% !TeX root = ../../surface-lagrangian-torus-fibration.tex
% 2026-11-23 / v1.10: Source-inclusive 正規形の共通像・正逆評価・全 orbit の定義域．
\section{共通の normal circle に適合する二段階平均}
\label{ke:staged-neck-section}

三重点の torus 平均を，neck 全体で共通に用いる normal circle の平均と，
end collar 上の残る circle の平均に分ける．
この順序によって，完成後の end tori が一つの actual Hamiltonian circle を含むことを
厳密に保証できる．単一 primitive を用いると normal moment の加法定数と
全 marked cycle の作用周期も同時に固定される．
平均化と Moser 法そのものの先行研究は
\cite[Section~3]{RuanHMinimal}を参照されたい．
以下の二段階構成と定量的適用の包括的な新規性確認は未完了である．

本節の円作用の周期はすべて $1$ とし，moment の符号は
\[
 \iota_V\omega=-d\mu
\]
に統一する．従って moment は一周の primitive 積分と同じ規格化である．
周期 $2\pi$ の生成場による通常の action 座標を用いる場合は，
生成場・moment・周期積分の規格化を同時に変更する．
また，以下の flow は指定する内側領域から出発して共通の外側領域に留まると仮定する．
作用・形式が一致する開集合を出る trajectory の一致は，局所的な ODE の一意性からは従わない．

\subsection{一つの円の平均と primitive の正確な輸送}

$U$ を複素曲面，$\omega=d\lambda$ を K\"ahler form とし，
$\rho_s$ を $U$ 上の正則円作用とする（固定点を許す）．
生成場を $V$，平均と chain homotopy を
\[
 \mathcal A\xi=\int_0^1\rho_s^*\xi\,ds,\qquad
 K\xi=\int_0^1(1-s)\rho_s^*(\iota_V\xi)\,ds
\]
と書く．Centered kernel $1/2-s$ に置き換えた場合も，
以下の恒等式と結論は同じである．

\begin{lemma}[円平均の exact Moser transport]
\label{ke:circle-exact-transport}
$\eta=\mathcal A\omega$，
$\omega_s=(1-s)\eta+s\omega$，
$\lambda_s=(1-s)\mathcal A\lambda+s\lambda$ とする．
$\beta=K\omega$ および
\[
 \iota_{Y_s}\omega_s=\beta,\qquad \phi_0=\mathrm{id}
\]
による flow が全時刻で定義される $\rho$-飽和した内側領域を取る．この領域では
\begin{equation}\label{ke:eq:circle-exact-transport}
 \phi_s^*\omega_s=\eta,\qquad
 \phi_s^*\lambda_s=\mathcal A\lambda+dq_s,
 \quad
 q_s=\int_0^s\phi_\tau^*
       \bigl(\lambda_\tau(Y_\tau)-K\lambda\bigr)\,d\tau .
\end{equation}
$\mu=(\mathcal A\lambda)(V)$ は $\eta$ の Hamiltonian moment である．
$\phi=\phi_1$ によって移した円作用
$\widehat\rho_s=\phi\circ\rho_s\circ\phi^{-1}$ は actual $\omega$ に関して Hamiltonian で，
その moment は
\begin{equation}\label{ke:eq:circle-actual-moment}
 \widehat\mu=\mu\circ\phi^{-1}.
\end{equation}
この moment の値は，対応する actual 円軌道上の $\lambda$ の一周積分に等しい．
\end{lemma}
\begin{proof}
Cartan の公式と部分積分により
\[
 dK+Kd=\mathcal A-I.
\]
従って
\[
 d\beta=\eta-\omega,\qquad
 \beta=\mathcal A\lambda-\lambda-d(K\lambda).
\]
正則作用による平均とその凸結合は K\"ahler form なので非退化である．
Moser の微分式から最初の pullback identity を得る．
Primitive については
\[
 \partial_s\lambda_s+\mathcal L_{Y_s}\lambda_s
 =\lambda-\mathcal A\lambda+\beta
     +d\bigl(\lambda_s(Y_s)\bigr)
 =d\bigl(\lambda_s(Y_s)-K\lambda\bigr)
\]
であり，flow で引き戻して積分すれば
\eqref{ke:eq:circle-exact-transport}の第二式となる．

$\mathcal A\lambda$ は円作用に不変なので
\[
 \iota_V\eta=-d\bigl((\mathcal A\lambda)(V)\bigr)=-d\mu.
\]
これを $\phi$ で移すと\eqref{ke:eq:circle-actual-moment}を得る．
周期 $1$ の軌道 $s\mapsto\rho_s(x)$ 上で
$(\mathcal A\lambda)(V)=\mu(x)$ は一定である．
Exact differential の周期は零なので，
$\int_{\phi(\rho_\bullet(x))}\lambda=\mu(x)$ である．
\end{proof}

\begin{lemma}[Actual circle の短さによる双方向 Moser buffer]
\label{ke:short-circle-actual-buffer}
上記の正則円作用と K\"ahler form に対し，作用不変な開集合 $U$ 全体で
$|V|_{g_\omega}\le\varepsilon$ と仮定する．
この補題の評価には $\omega$ の exactness，作用不変な背景計量，
tensor の二側比較は不要である．
$w(a)=1-a$ または $w(a)=1/2-a$ とし，
\[
 \beta=\int_0^1 w(a)\rho_a^*(\iota_V\omega)\,da,
 \quad \eta=\mathcal A\omega,
 \quad \omega_s=(1-s)\eta+s\omega,
 \quad \iota_{Y_s}\omega_s=\beta
\]
と置く．$c_w=(\pi/2)\|w\|_{L^2(0,1)}$ とすれば，
\begin{equation}\label{ke:eq:circle-actual-speed}
 |Y_s|_{g_\omega}
 \le\frac{\varepsilon\|w\|_{L^2(0,1)}}{2\sqrt{s(1-s)}}
 \qquad (0<s<1).
\end{equation}
従って，任意の二つの時刻を結ぶ存在する trajectory の actual 長さは
\begin{equation}\label{ke:eq:circle-actual-length}
 \operatorname{Length}_{g_\omega}(\phi_{s,s_0}(x))
 \le c_w\varepsilon,
 \qquad
 c_w=\begin{cases}
 \pi/(2\sqrt3),&w(a)=1-a,\\
 \pi/(4\sqrt3),&w(a)=1/2-a
 \end{cases}
\end{equation}
である．ここで時間を逆向きに辿る場合も同じ評価を用いる．

$K\subset U$ を compact な $\rho$-飽和集合とする．
$K\subset D\Subset U$ なる開集合 $D$ が存在し，
\begin{equation}\label{ke:eq:circle-actual-buffer}
 \delta:=d_{g_\omega}^{U}(K,\partial D)>c_w\varepsilon
\end{equation}
を満たすとする．$d_{g_\omega}^{U}$ は $U$ 内の経路による actual 距離である．
このとき $K$ のある $\rho$-飽和開近傍からの forward flow は全時刻で存在する．
また，時刻 $1$ に $K$ の点を終点として指定する backward flow も時刻 $0$ まで存在する．
従って最大の time-one flow の定義域と像はともに $K$ を含む．
各円軌道 $\mathcal O\subset K$ に対して
\begin{equation}\label{ke:eq:circle-actual-diameter}
 \operatorname{diam}_{d_{g_\omega}^{U}}\phi_{1,0}(\mathcal O)
 \le (2c_w+1/2)\varepsilon.
\end{equation}
ここで直径は ambient 距離による．
像の円自身の intrinsic diameter，flow の $C^1$ 評価，
元の neck base projection の保存は結論に含めない．
\end{lemma}
\begin{proof}
$g_a=\rho_a^*g_\omega$ と置く．正則性により
$g_\eta=\int_0^1g_a\,da$ であり，$d\rho_a(V)=V$ と仮定から
$g_a(V,V)\le\varepsilon^2$ である．任意の接ベクトル $X$ に対し
\begin{align*}
 |\beta(X)|
 &\le\varepsilon\int_0^1 |w(a)|\,|X|_{g_a}\,da\\
 &\le\varepsilon\|w\|_{L^2(0,1)}
       \left(\int_0^1g_a(X,X)\,da\right)^{1/2}
 =b|X|_{g_\eta},\qquad b=\varepsilon\|w\|_{L^2(0,1)}.
\end{align*}
共通複素構造を $J$ とする．Moser 方程式から
$g_{\omega_s}(Y_s,Y_s)=\beta(JY_s)$ なので，
$r=|Y_s|_{g_\eta}$，$q=|Y_s|_{g_\omega}$ と置けば
\[
 (1-s)r^2+s q^2\le br.
\]
平方完成によって
\[
 s q^2\le br-(1-s)r^2\le\frac{b^2}{4(1-s)}
\]
となり，\eqref{ke:eq:circle-actual-speed}を得る．
$\int_0^1(s(1-s))^{-1/2}\,ds=\pi$ を用いて
\eqref{ke:eq:circle-actual-length}が従う．
端点で表示した上界は発散するが可積分である．
$Y_s$ 自身は K\"ahler 形式 $\omega_s$ の非退化性により
$[0,1]\times U$ 上で滑らかに定義されている．

任意の時刻に $K$ から出発する trajectory が初めて $D$ の境界へ達すれば，
そこまでの長さは $\delta$ 以上になる．
これは\eqref{ke:eq:circle-actual-length}に反する．
従って trajectory は $\overline D\Subset U$ に留まり，
compact 集合上の滑らかな時間依存 ODE の延長定理によって
順方向・逆方向とも $[0,1]$ 全体へ延びる．
Strict な余裕と $K$ の compact 性により，この結論は $K$ の開近傍にも成立する．
$K$ と作用群が compact なので，その近傍に含まれる
$\rho$-飽和開近傍を選べる．Backward flow の時刻 $0$ の点は forward flow の定義域に属するから，
最大 flow の像も $K$ を含む．

元の円軌道の任意の二点は高々半周で結べ，その長さは $\varepsilon/2$ 以下である．
両端にそれぞれ長さ $c_w\varepsilon$ 以下の Moser trajectory を付ければ
\eqref{ke:eq:circle-actual-diameter}となる．
\end{proof}

\begin{corollary}[同じ compact 領域を覆う actual 円作用]
\label{ke:short-circle-image-buffer}
補題\ref{ke:short-circle-actual-buffer}の仮定に加えて
\begin{equation}\label{ke:eq:circle-image-buffer}
 \delta>(2c_w+1/2)\varepsilon
\end{equation}
とする．このとき time-one Moser map $\phi$ の定義域には
$\rho$-飽和した開集合 $O\supset K$ を取ることができ，その像も $K$ を含む：
\[
 K\subset O,\qquad K\subset\phi(O).
\]
従って $\widehat\rho=\phi\rho\phi^{-1}$ は $K$ の開近傍上の円作用となる．
さらに $\omega=d\lambda$ の場合は補題\ref{ke:circle-exact-transport}の
actual Hamiltonian action と moment の結論がその近傍で成立する．
\end{corollary}
\begin{proof}
$F$ を，時刻 $0$ からの trajectory が全時刻で $D$ 内に留まる点の集合とする．
これは開集合であり，time-one map は $F$ 上の滑らかな埋込である．
その最大の $\rho$-飽和部分
\[
 O=\{y\in F:\rho_a(y)\in F\text{ for every }a\in\mathbb R/\mathbb Z\}
\]
も開である．実際，ある軌道が開集合 $F$ に含まれれば，軌道の compact 性と
作用の連続性によりその近くの全軌道も $F$ に含まれる．
補題の forward 結論と $K$ の飽和性から $K\subset O$ である．

$x\in K$ を取り，backward flow による時刻 $0$ の点を $y$ とする．
$d_{g_\omega}^{U}(x,y)\le c_w\varepsilon$ である．
任意の $z=\rho_a(y)$ は $y$ から高々半周で結べるので
\[
 d_{g_\omega}^{U}(x,z)\le(c_w+1/2)\varepsilon.
\]
$y$ から $z$ への短い軌道弧も $D$ を出られない．
もし出れば，$x$ から $\partial D$ までの経路の長さが
$(c_w+1/2)\varepsilon<\delta$ となるからである．
$z$ からの forward Moser trajectory が $D$ を出るなら，
先ほどの $x$ から $z$ までの経路と合わせて
$x$ から $\partial D$ まで長さ $(2c_w+1/2)\varepsilon$ 以下の経路ができる．
これは\eqref{ke:eq:circle-image-buffer}に反する．
従って全ての $z=\rho_a(y)$ は $F$ に属し，$y\in O$ となる．
$x=\phi(y)$ だから $K\subset\phi(O)$ が従う．
最後の主張は既存の exact transport identity を適用すればよい．
\end{proof}

\begin{remark}[Neck への適用範囲]
上記の二つの結果は，actual normal generator の一様な短さと
actual 距離の buffer から normal averaging を実行する条件である．
Rank-one 評価と volume 比較から短さを得る場合も，
source-free neck 全体でこれらの入力が一様であることは別の評価を要する．
固定残余基底上では定理\ref{ke:fixed-base-neck-moser}で検証し，
source を含む固定段階への拡張は補題\ref{ke:fixed-stage-source-neck}で扱う．
複数の triple end における同一作用，reduced leaves，markings の同時整合，
source・branch・core への接着や全曲面の GH 収束は，この補題だけの結論ではない．
それぞれ系\ref{ke:neck-corner-saturated-union}，
定理\ref{ke:actual-global-fibration}，および
定理\ref{ke:global-ambient-lower}，\ref{ke:global-ambient-upper}で証明する．
\end{remark}

\subsection{固定残余領域上の法線作用の平均エネルギー}

\begin{proposition}[法線方向の積分評価と短い円軌道]
\label{ke:neck-integrated-normal-energy}
Source と三重点から離れた固定座標で $uv=t$，$z\in D$ とし，
$|u|,|v|<r$ に対して $s=-\log(|u|^2/r^2)\in(0,M_t)$ と置く．
ここで $M_t=-\log(|t|^2/r^4)$ である．
固定した $Q\Subset D$ 上で，actual negative KE form $\omega_t$ の体積が
annulus--disc complete KE 積計量の体積と $a_0\le s\le M_t-a_0$ で
一様に二側比較可能とする（$a_0>0$ は固定）．
$V=2\pi\partial_{\arg u}$，$\rho$ をその周期 $1$ の作用とし，
\[
 E_t(s,z)=\int_0^1 |V|_{g_t}^2(\rho_a(u,z))\,da
\]
と置く．$2a_0\le m\le M_t/2$ ならば，全 $z\in Q$ について
\begin{equation}\label{ke:eq:normal-energy-mass}
 \int_m^{M_t-m} E_t(s,z)\,ds\le C/m .
\end{equation}
定数 $C$ は $t,m,z$ に依存しない．任意の $A>0$ について
$[m,M_t-m]$ 内で $E_t>A/m^2$ となる深度の Lebesgue 測度は
$Cm/A$ 以下である．それ以外の深度では actual 円軌道の長さは
$\sqrt A/m$ 以下である．
\end{proposition}

\begin{proof}
$\Omega=du/u\wedge dz$，$F_t=\log(\omega_t^2/|\Omega|^2)$ とし，
$\bar F_t$ を $u$ の角度について平均する．$|\Omega|^2$ は
$\Omega$ の定める正の体積形式で，その定数規格化は以下に影響しない．
$d^c=\tfrac{\sqrt{-1}}2(\bar\partial-\partial)$ の規格化では
$dd^cF_t=\omega_t$ である．$d(s)=\min(s,M_t-s)$ とすると，
積計量の体積密度は $z\in Q$ 上で
$(\pi/M_t)^2\csc^2(\pi s/M_t)\asymp d(s)^{-2}$ に比較可能だから
\[
 \bigl|\bar F_t(s,z)+2\log d(s)\bigr|\le C_0 .
\]
局所対数座標 $\xi=\log u$ で $s=-2\Re\xi+2\log r$ である．
$\mathcal A\omega_t=dd^c\bar F_t$ の $\xi\bar\xi$ 係数から
$E_t=8\pi^2\partial_s^2\bar F_t$ を得る．特に $\bar F_t$ は $s$ の凸関数である．

$s\in[m,M_t-m]$ で $s\pm m/2$ を取ると，三点の $d$ の値の比は
$1/2$ と $3/2$ の間にある（$d(s)$ を分母とする）．
従って前進・後退の差商は絶対値 $C_1/m$ 以下であり，
凸性から $|\partial_s\bar F_t|\le C_1/m$ を得る．
両端における導関数の差を取れば
\eqref{ke:eq:normal-energy-mass} が従う．
測度評価は Markov の不等式，長さ評価は
$\int_0^1 |V|_{g_t}(\rho_a x)\,da\le E_t(s,z)^{1/2}$ による．
\end{proof}

上の体積仮定は，一般位置全平面退化の固定 source-free pair 座標では，
関連 KE 主稿の \texttt{neck:prop:volume-envelope} から従う．
例えば $m=R_t$，$A=R_t$ とすれば，各固定残余点で中央 strip のうち
深度の測度 $O(1)$ の例外を除いて円軌道長は $O(R_t^{-1/2})$ である．
この例外集合は残余点と $t$ に依存する．

\begin{corollary}[鋭い係数を持つ同一残余点での法線接続上界]
\label{ke:neck-axial-upper}
命題\ref{ke:neck-integrated-normal-energy}の座標で $z\in Q$ を固定する．
$2a_0\le a\le b\le M_t-2a_0$ とする．ある一定の角度で深度 $a$ から $b$ へ
進む動径弧の actual 長さは
\[
 \int_a^b \frac{ds}{d(s)}+C
\]
以下である．従って両端円の任意の点を，同じ $z$ を保つ経路で
この上界に $\tfrac12(\sqrt{E_t(a,z)}+\sqrt{E_t(b,z)})$ を加えた
長さ以下で結べる．
\end{corollary}

\begin{proof}
まず $b\le M_t/2$ とする．$\bar F_t=-2\log s+O(1)$ で，
命題の差商評価に各点で $m=d(s)=s$ を用いると $|s\partial_s\bar F_t|\le C_1$
である．部分積分して
\[
 \int_a^b s\partial_s^2\bar F_t\,ds
 =[s\partial_s\bar F_t]_a^b-[\bar F_t]_a^b
 \le 2\log(b/a)+C_2 .
\]
$JV=4\pi\partial_s$ と $J$ 不変性から，
$g_t(\partial_s,\partial_s)$ の角度平均は
$E_t/(16\pi^2)=\tfrac12\partial_s^2\bar F_t$ である．
従って Cauchy--Schwarz により動径弧長の角度平均は
\[
 \int_a^b \sqrt{\tfrac12\partial_s^2\bar F_t}\,ds
 \le \sqrt{\log(b/a)\bigl(\log(b/a)+C_2/2\bigr)}
 \le \log(b/a)+C_2/4 .
\]
上側半区間では $s$ を $M_t-s$ に置き換える．
中点を横切る場合は角度平均を取った後で両半区間の上界を加えるので，
全区間に共通する一つの角度を選べる．
両端では円の短い方の弧を付加する．
\end{proof}

この接続上界は正規化前の距離に対し加法誤差 $O(1)$ を持つ．
両端の円エネルギーも有界な場合は，正規化した法線接続費用が
$R_t^{-1}\int_a^b ds/d(s)+O(R_t^{-1})$ 以下となる．
この評価だけでは異なる残余点間の portal 接続費用や，全曲面を通る経路の下界は得られない．
前者は命題\ref{ke:fixed-base-moment-connectors}と
系\ref{ke:source-inclusive-same-moment-collapse}，
後者は定理\ref{ke:global-ambient-lower}で別に評価する．

\begin{proposition}[固定残余基底上の moving deep moment strip]
\label{ke:neck-deep-moment-strip}
関連 KE 主稿の \texttt{pair:prop:bundle-conic} の固定 source-free compact
基底 $Q\Subset D^\circ$ 上で，双対 Hermitian norms を用いて
$uv=t$，$s=-\log(\|v\|^2/r^2)$，$M_t=-\log(|t|^2/r^4)$ とする．
$V$ は $v$ を正方向に回す周期 $1$ の生成場とし，
$\mu_t=(\mathcal A\lambda_t)(V)$ は同稿 \texttt{pair:cor:global-primitive}
の canonical moment とする．
定数 $m_0,c,C>0$ と $\vartheta\in(0,1/2)$ を固定できて，
$m_0\le m\le\vartheta M_t$ のすべての深度について，全 $Q$ 上で
\[
 \frac c m\le\mu_t\big|_{s=m}\le\frac C m,
 \qquad
 -\frac C m\le\mu_t\big|_{s=M_t-m}\le-\frac c m
\]
が成り立つ．従って任意の $m=m_t$ がこの範囲にあれば，
共通 interval $I_t=[-c/(2m_t),c/(2m_t)]$ の各 level は
$m_t<s<M_t-m_t$ 内で各 annulus に一つの円を持つ．
その商は $I_t\times Q$ であり，$Q$ が境界付き領域なら境界も保つ．
特に $\mu_t=0$ の level は両側深度が少なくとも $\vartheta M_t$
である領域に入る．
\end{proposition}

\begin{proof}
局所正則 frame で $\Omega=dv/v\wedge dz$ をとる．
$\bar F_t=\mathcal A\log(\omega_t^2/|\Omega|^2)$ とすると，
定数倍および base-only factor を除いた primitive の規格化から
$\mu_t=-2\pi\partial_s\bar F_t$ を得る．
Hermitian norm による $s$ と正則 frame の座標による深度の差は
固定 compact 上で有界な base-only 関数であり，$\partial_s$ は変わらない．
有限個の少し大きい積座標に命題\ref{ke:neck-integrated-normal-energy}で使った
体積包絡を適用して，$d(s)=\min(s,M_t-s)$ とすれば
ある固定 $a_0>0$ について
\[
 |\bar F_t(s,z)+2\log d(s)|\le C_0
 \qquad (a_0\le s\le M_t-a_0)
\]
を $Q$ 全体に一様に得る．Base coordinate の変更による加法項は
有限個の compact patches 上で有界であり，差商には現れない．

$2\log\Lambda>2C_0+1$ となる固定 $\Lambda>1$ を選び，
$\vartheta=1/(2\Lambda)$ とする．$2a_0\le m\le\vartheta M_t$ なら，
$m$ と $\Lambda m$ は左側半区間にある．凸性により
\[
 \partial_s\bar F_t(m,z)
 \le\frac{\bar F_t(\Lambda m,z)-\bar F_t(m,z)}{(\Lambda-1)m}
 \le -\frac{1}{(\Lambda-1)m}.
\]
逆向きの大きさの上界は命題\ref{ke:neck-integrated-normal-energy}の
差商評価 $|\partial_s\bar F_t|\le C_1/m$ による．
右側は $M_t-s$ を使って同様に求める．
$\partial^2_s\bar F_t>0$ なので $\mu_t$ は各 annulus の深度に対し
厳密に減少する．中間値定理と陰関数定理で，$I_t$ 全体にわたる
一意な深度断面が滑らかに貼り合わさる．半径を正に変えても
円束の標識は保たれるので商は $I_t\times Q$ になる．
最後は $m=\vartheta M_t$ を選ぶ．
\end{proof}

この interval は normal 深度を動かしつつ，すべての固定残余点で
同じ canonical moment を使うためのものである．
既存の固定 end の moment 漸近を moving cut に代入する必要はない．
レベル上で点ごとの円エネルギーを抑える評価と actual Moser buffer は，
定理\ref{ke:fixed-base-neck-moser}で別に与える．
Source を含む固定段階と三重点へ動く collars の同時選択には，
補題\ref{ke:fixed-stage-source-neck}と定理\ref{ke:neck-corner-diagonal-matching}を用いる．

\begin{corollary}[軌道平均エネルギーを用いる Moser 速度評価]
\label{ke:neck-averaged-speed}
補題\ref{ke:short-circle-actual-buffer} の Moser 方程式において，
点 $x$ で
$E(x)=\int_0^1 |V|_{g_\omega}^2(\rho_a x)\,da$ とすると，
\[
 |Y_\tau(x)|_{g_\omega}
 \le \frac{\|w\|_{L^\infty}\sqrt{E(x)}}{2\sqrt{\tau(1-\tau)}} ,
 \qquad 0<\tau<1 .
\]
作用不変な開集合上で $E\le e^2$ が一様に成り立てば，
補題\ref{ke:short-circle-actual-buffer} と
系\ref{ke:short-circle-image-buffer} の長さ・buffer の結論は
$\varepsilon=e$，$c_w=(\pi/2)\|w\|_{L^\infty}$ で成り立つ．
\end{corollary}

\begin{proof}
$g_a=\rho_a^*g_\omega$ とすると
\[
 |\beta(X)|\le
 \left(\int_0^1 |w(a)|^2 g_a(V,V)\,da\right)^{1/2}
 |X|_{g_\eta}\le \|w\|_{L^\infty}\sqrt{E(x)}\,|X|_{g_\eta}.
\]
この定数を補題\ref{ke:short-circle-actual-buffer} の平方完成に用いる．
また $E\le e^2$ 上で一周の長さが $e$ 以下なので，
二点間の短い方の軌道弧の長さは $e/2$ 以下である．
これを使って両向きの軌道追跡と飽和近傍の証明をそのまま適用できる．
\end{proof}

積分評価\eqref{ke:eq:normal-energy-mass}は所与の深度での点ごとの上界や，
内側集合から小エネルギー領域の境界までの actual 距離下界を与えない．
例えば残余方向を固定しても凸関数の二階微分は小さい区間に集中し得る．
被覆を得るには，所要の開 collar 上の一様な $E$ の上界と actual buffer，
または軌道ごとに同じ被覆を保証する変位評価が必要となる．
本稿では定理\ref{ke:fixed-base-neck-moser}で前者を固定残余基底上に構成し，
補題\ref{ke:fixed-stage-source-neck}で source を含む場合へ拡張する．

\begin{corollary}[縮小する actual buffer からの Moser 被覆]
\label{ke:neck-shrinking-buffer}
上記の annulus--disc 座標で，compact な飽和集合 $K_t$ を含む飽和開集合
$U_t\subset\{d(s)>m_t\}$ を取る．$g_t$ に関する全曲面の距離で
\[
 \delta_t=\min\bigl(1,\operatorname{dist}_{g_t}(K_t,X_t\setminus U_t)\bigr)
\]
と置く．$m_t\delta_t^2\to\infty$ ならば，十分小さい $t$ で
$K_t$ を source と像の双方に含む飽和 source を持つ normal Moser map が存在する．
その円軌道像の ambient 直径は $O((m_t\delta_t)^{-1})$ である．
\end{corollary}

\begin{proof}
$N_t=\{\operatorname{dist}_{g_t}(\cdot,K_t)<\delta_t/2\}$ と置く．
各 $q\in N_t$ の actual $\delta_t/4$-ball の閉包は $U_t$ に含まれる．
残余射影への localized Chern--Lu 評価を半径 $\delta_t/4$ で用いると，
中心で $p^*h_D\le C\delta_t^{-2}\omega_t$ を得る．
この半径依存性は，距離 cutoff の勾配二乗と Laplacian の下界を
$O(\delta_t^{-2})$ で制御する Ricci 下界による距離比較から従う．
関連 KE 主稿の \texttt{neck:lem:rank-one-circle} と体積上界により
$N_t$ 上で $|V|_{g_t}\le\varepsilon_t=C'/(m_t\delta_t)$ となる．

$D_t=\{\operatorname{dist}_{g_t}(\cdot,K_t)<\delta_t/4\}$ とする．
$\overline D_t$ から任意の位相まで周期の半分以下で円軌道を辿るとき，
$N_t$ の最初の退出までの長さは $\varepsilon_t/2$ 以下である．
これは $\delta_t/4$ より短いので，$\overline D_t$ を通る全軌道が $N_t$ に留まる．
従って $N_t$ に含まれる最大の飽和開集合に $\overline D_t$ が含まれ，
そこで同じ短円上界を使える．$K_t$ から $\partial D_t$ までの actual 距離は
$\delta_t/4$ 以上であり，$m_t\delta_t^2\to\infty$ により
$(2c_w+1/2)\varepsilon_t$ を上回る．系\ref{ke:short-circle-image-buffer}と
補題\ref{ke:short-circle-actual-buffer}の直径評価を適用する．
\end{proof}

\subsection{固定残余領域での actual tensor 比較と Moser 被覆}

\begin{lemma}[大域 capped density を用いた射影計量の下界]
\label{ke:capped-projective-lower}
固定次数 $d\ge5$ の一般位置全平面退化を考え，
$\Omega_{t,d}^{\rm low}$ と $h_{t,d}^{\rm low}$ を
関連 KE 主稿の \texttt{coarse:cor:degree-d-volume-lower} で定める．
$\varphi_t=\log(\omega_t^2/\Omega_{t,d}^{\rm low})$ とすると，
ある固定 $A,C>0$ に対し
\[
 \operatorname{tr}_{\omega_t}(\omega_{\mathrm{FS}}|_{X_t})
 \le C\exp(A\varphi_t)
\]
が全 $X_t$ で成り立つ．従って $\varphi_t$ が一様に上から有界な領域では
$g_t\ge c g_{\mathrm{FS}}|_{X_t}$ となる．
\end{lemma}

\begin{proof}
KE 側の capped density の証明で $A_0$ を十分大きく固定して，
$h_{t,d}^{\rm low}\ge c_0\omega_{\mathrm{FS}}|_{X_t}$ を得る．
実際，背景は $(d-4)\omega_{\mathrm{FS}}|_{X_t}$ で，
すべての cap の負部分と frame errors の和は
その $O_d(A_0^{-1})$ 倍に抑えられ，残りは半正である．
同じ大域体積下界から $\varphi_t\ge-C_0$ が全 $X_t$ で成り立つ．
$T_t=\operatorname{tr}_{\omega_t}(\omega_{\mathrm{FS}}|_{X_t})$ とすると，
正則包含 $X_t\to\PP^3$ の Chern--Lu 不等式は
$\Delta_{\omega_t}\log T_t\ge-C_1-C_2 T_t$ を与える．
ここでは射影空間の固定した curvature upper bound と
$\operatorname{Ric}(\omega_t)=-\omega_t$ を用いる．
$dd^c\varphi_t=\omega_t-h_{t,d}^{\rm low}$ なので
\[
 \Delta_{\omega_t}(\log T_t-A\varphi_t)
 \ge (Ac_0-C_2)T_t-(C_1+2A).
\]
$Ac_0>C_2$ と取り，compact $X_t$ 上の最大点で $T_t$ を上から抑える．
そこでの $\varphi_t\ge-C_0$ と合わせ，
$\log T_t-A\varphi_t\le C'$ を全域に得る．
\end{proof}

\begin{theorem}[固定 source-free base の moving strip での actual Moser 被覆]
\label{ke:fixed-base-neck-moser}
関連 KE 主稿の \texttt{pair:prop:bundle-conic} の正則束座標を用い，
source と三重点を避けた固定した $Q\Subset B$ 上で
$uv=t$，$s=-\log(\|v\|^2/r^2)$，
$M_t=-\log(|t|^2/r^4)$，$d(s)=\min(s,M_t-s)$ とする．
$Q$ は compact とする．固定近傍に依存する有限な $m_Q$ がある．
$|t|$ を十分小さくし，$m_Q\le m_t\le M_t/2$ に対して
$K_t=\{p\in Q,\ d(s)\ge m_t\}$ とする．
$K_t$ を source と像の双方に含む normal Moser map が存在し，
その軌道像の actual ambient 直径は $O(m_t^{-1})$ である．
さらに $K_t$ を含む深部開 strip 上で，局所積座標の complete
annulus--disc 計量 $g_{\rm prod}$ と $g_t$ は一様に二側比較可能である．
定数は $m_t,t$ に依存しない．従って固定した十分大きい cut にも使え，
$m_t\to\infty$ と選べば軌道の変位は零へ収束する．
\end{theorem}

\begin{proof}
まず $\overline Q$ の少し大きい固定残余近傍を有限個の局所積座標で覆う．
各 patch で outer annulus--disc を固定し，その内部に
三重に入れ子になった annulus--disc を，正の固定な normal radii と residual radii
で選ぶ．最も外側の入れ子領域（outer chart には compact に含まれる）では
outer chart の体積 envelope と capped density が一様に比較可能である．
これは \texttt{neck:prop:volume-envelope} の証明にある
ordinary-pair capped density の比較である．同命題の体積上界から
そこでは $\varphi_t\le C$ であり，補題\ref{ke:capped-projective-lower}
によって $g_t\ge c g_{\mathrm{FS}}$ を得る．

内側の入れ子領域から中間領域の外へ出る経路は，その最初の退出までに，
残余座標又は $u,v$ の一方を固定した正の幅だけ変化させる．
これらの座標は固定な ambient 座標上で有界な微分を持つ．
実際，正則束座標は $t$ に依存しない ambient 変換で作られ，
コンパクトな座標範囲では変換とその逆の微分が一様に有界である．
よって経路は比較が成立する部分だけで一様に正の actual 長さを持つ．
従って内側の各点の一様な半径の actual ball の閉包は中間領域に含まれる．

この actual ball 上で outer annulus--disc の complete product への
正則恒等写像に localized Chern--Lu 評価
（関連 KE 主稿 \texttt{coarse:lem:localized-chern-lu}）を適用する．
曲率 $-1$ の二曲面の積の holomorphic sectional curvature は
$-1/2$ 以下なので，中心で $g_t\ge c g_{\rm prod}$ を得る．
既存の体積上界と合わせれば複素次元 $2$ の行列の固有値比較により
$g_t\le C g_{\rm prod}$ も従う．有限被覆で定数を一様にできる．
固定 norm による $s$ は局所積座標の深度と有界な base-only 項だけ異なり，
その項の微分も一様に有界である．従って十分深い領域上で
\[
 |V|_{g_t}\le C/d(s),\qquad
 |dp|_{g_t}\le C,\qquad
 |d\log d(s)|_{g_t}\le C
\]
となる．最後の不等式は $s=M_t/2$ の継目では片側微分，
又は経路に沿う Lipschitz 評価として読む．$|dp|$ の標的には固定基底計量を用いる．

$Q\Subset B'\Subset B$ を固定して，Moser の外側領域を
$U_t=\{p\in B',\ d(s)>m_t/2\}$ とする．
$K_t$ から $U_t$ の境界へ向かう任意の経路は，最初の退出以前に
基底で固定した正の距離を移動するか，$\log d(s)$ を $\log 2$ 以上変える．
従って actual buffer は正の定数以上である．ここで局所比較のための有限被覆は
$\overline{B'}$ 上に選んでいる．同じ議論で $B'$ と $m_t/2$ を少し内側へ
縮めた比較的 compact な領域を $U_t$ 内に取れる．
$U_t$ 上の $|V|_{g_t}\le C_Q/m_t$ を用いる．上で得た正の buffer を
$\beta_Q$ とし，$m_Q$ を
$(2c_w+1/2)C_Q/m_Q<\beta_Q$ となるまで大きく取れば，
全ての $m_t\ge m_Q$ について系\ref{ke:short-circle-image-buffer}を適用できる．
この条件も有限被覆の tensor 比較も $m_t\to\infty$ 自体を必要としない．
従って固定 cut を含めて求める Moser map を得る．
直径は補題\ref{ke:short-circle-actual-buffer}の評価による．
\end{proof}

定理\ref{ke:fixed-base-neck-moser} は residual base が固定した場合を扱う．
Source や三重点へ近付く residual collars では，射影計量で測った入れ子領域間の幅や
capped density と complete envelopes の比較を別に用意する．
Source は定理\ref{ke:projected-source-neck}，三重点の全 end との整合は
定理\ref{ke:neck-corner-diagonal-matching}と系\ref{ke:neck-corner-saturated-union}で扱う．
固定 compact の結論と命題\ref{ke:neck-deep-moment-strip} を合わせると，
同じ固定基底上で moving canonical moment strip を actual form に移せる．
複数 end の接着では，これに加えて上記の正則束作用と triple 作用との
開 overlap 上での整合を用いる．

\begin{proposition}[同一 normal moment 上の残余方向を含む有界な接続]
\label{ke:fixed-base-moment-connectors}
命題\ref{ke:neck-deep-moment-strip}と定理\ref{ke:fixed-base-neck-moser}の
固定 source-free 基底で，$P$ を連結な compact 境界付き曲面とし，
その少し大きい固定近傍上で両結果を用いる．
固定近傍に依存する十分大きい有限な $m_P$ を取り，
$m_P\le m_t\le\vartheta M_t$ とし，
$I_t=[-c/(2m_t),c/(2m_t)]$ を共通 moment interval とする．
円平均後の canonical moment を $\mu_t$，同じ normal Moser map を
$\phi_t$ とし，$a\in I_t$ に対して
\[
 E_t(a)=\{p\in P,\ m_t\le s\le M_t-m_t,\ \mu_t=a\}
\]
と置く．$t,a,m_t$ に依存しない $C_P$ が存在して
\begin{equation}\label{ke:eq:fixed-base-moment-connectors}
 \operatorname{diam}_{d_{g_t}}E_t(a)\le C_P,
 \qquad
 \operatorname{diam}_{d_{g_t}}\phi_t(E_t(a))
 \le C_P+O(m_t^{-1})
\end{equation}
となる．従って actual normal moment の同じ level に属する
この neck 内の全 fibres の合併は，正規化距離 $R_t^{-1}d_{g_t}$ で
一様に $O(R_t^{-1})$ の直径を持つ．
接続経路は moment level や個々の fibre の中に留まる必要はなく，
結論は ambient diameter についてのものである．
\end{proposition}

\begin{proof}
$d(s)=\min(s,M_t-s)$ とする．
命題\ref{ke:neck-deep-moment-strip}の証明から，定数を共通に取り
\[
 |\mu_t(s,p)|\le C_0/d(s),\qquad
 c_0/s\le\mu_t(s,p)\le C_0/s
 \quad(m_0\le s\le\vartheta M_t)
\]
を得る．右端では $s$ を $M_t-s$ に置き換え，符号を逆にする．
各 $p$ で $s\mapsto\mu_t(s,p)$ は厳密に減少する．
$A=2C_0/\vartheta$ と取ると，$a>A/M_t$ の level はすべて左端側にあり，
\[
 c_0/a\le s(a,p)\le C_0/a.
\]
$a<-A/M_t$ なら $M_t-s(a,p)$ が同様に比較できる．
残る $|a|\le A/M_t$ では，左端側にある点に上の下界を適用すると
$s(a,p)\ge c_0M_t/A$ となる．右端側も同様なので，
ある固定 $\kappa>0$ に対して
\[
 \kappa M_t\le s(a,p)\le(1-\kappa)M_t
\]
を得る．以上の三場合のいずれでも，任意の二つの残余点について
\begin{equation}\label{ke:eq:moment-depth-connector}
 \left|\int_{s(a,p)}^{s(a,q)}\frac{ds}{d(s)}\right|\le C_1
\end{equation}
が一様に成り立つ．この議論は $s(a,p)$ の残余微分を評価しない．
用いた $c_0,C_0,\vartheta$ と深度の下限 $m_0$ は固定近傍だけに依存するので，
$m_P$ を $m_0$ と前定理の Moser 下限以上に取れば，固定した有限 cut にも
同じ証明が適用される．特に固定 stage を先に選び，その後 $t\to0$ とする用途に使える．

固定した円束の smooth connection と $P$ 上の固定補助計量を取る．
連結 compact 曲面 $P$ の任意の二点は，この補助計量で長さが一様に有界な
区分的 smooth な経路で結べる．定理\ref{ke:fixed-base-neck-moser}の
二側比較を有限個の局所積座標に適用すると，深度 $s$ を固定して
この経路を水平に持ち上げた actual 長さも一様に有界である．
実際，局所正則 frame と固定 Hermitian norm の深度の差，および
connection form は固定基底上で一様有界であり，normal 方向の寄与は
$d(s)^{-1}$ 倍される．同じ二側比較から，固定残余点・固定角度の
動径弧の長さは $C\int ds/d(s)$ 以下である．

そこで $E_t(a)$ の二点を，第一点の深度を保つ水平経路，
第二点の残余点での動径経路，最後の短い円弧によって結ぶ．
動径費用は\eqref{ke:eq:moment-depth-connector}，
円弧の費用は $|V|_{g_t}=O(m_t^{-1})$ から有界である．
全経路は深部 neck 内に取り，最初の不等式を得る．
最後に両端へ $\phi_t$ の trajectory を付ける．
各 trajectory の actual 長さは $O(m_t^{-1})$ なので第二式が従う．
これは同一 $a$ 上の任意の二点に一様であり，reduced foliation の
葉の選択や葉内での長さに依存しない．
\end{proof}

\begin{corollary}[異なる moment を含む固定基底の中央帯]
\label{ke:fixed-base-central-band}
命題\ref{ke:fixed-base-moment-connectors}と同じ固定連結 source-free 基底 $P$，
canonical moment $\mu_t$，normal Moser map $\phi_t$ を用いる．
$A>0$ を固定し，前命題の定数を用いて
\[
 m_P\le m_t\le\min\{\vartheta,c/(2A)\}M_t
\]
とする．これは固定した十分大きい cut を含み，$m_t=o(M_t)$ の選択でも満たされる．
十分小さい $t$ について
\[
 E_{t,A}=\{p\in P,\ m_t\le s\le M_t-m_t,\ |\mu_t|\le A/M_t\}
\]
と置くと，$t,m_t$ に依存しない $C_{P,A}$ が存在して
\[
 \operatorname{diam}_{d_{g_t}}E_{t,A}\le C_{P,A},\qquad
 \operatorname{diam}_{d_{g_t}}\phi_t(E_{t,A})
 \le C_{P,A}+O(m_t^{-1})
\]
が成り立つ．異なる normal moment に属する点も同時に含み，
正規化した ambient diameter は $O_{P,A}(R_t^{-1})$ である．
\end{corollary}

\begin{proof}
$m_t\ge m_P$ により，前命題の片側 moment 評価に必要な固定最小深度を超える．
また仮定の上限から $m_t\le\vartheta M_t$ および $A/M_t\le c/(2m_t)$ となるので，
中央帯の moment 値は同じ共通 interval に入る．
$s\le\vartheta M_t$ での $\mu_t\ge c_0/s$ と
$|\mu_t|\le A/M_t$ を合わせれば $s\ge c_0M_t/A$ を得る．
右側では符号を逆にした評価を用いる．従って
\[
 \kappa_A=\min\{\vartheta,c_0/A,1/4\}>0
\]
と置けば，中央帯の全点で
$\kappa_A M_t\le s\le(1-\kappa_A)M_t$ である．
特に，異なる moment に属する任意の二点の深度 $s_1,s_2$ に対しても
\[
 \left|\int_{s_1}^{s_2}\frac{ds}{\min(s,M_t-s)}\right|
 \le 2\log\frac{1}{2\kappa_A}
\]
となる．前命題の固定深度での水平 lift，終点の残余点での動径弧，
最後の phase 弧を用いれば，この二点を一様な長さで結べる．
経路は深部 neck 内に留まるが，中央帯や moment level 内に留まる必要はない．
両端に長さ $O(m_t^{-1})$ の Moser trajectories を加えれば像の評価も従う．
\end{proof}

以上の命題と系は，固定した source-free 基底上で異なる残余点を結ぶ portal 費用を扱う．
Source の pinched fibre を含む全共通 interval の同一 moment 集合は，
系\ref{ke:source-inclusive-same-moment-collapse}で評価する．
三重点へ動く残余 collars は固定段階の exhaustion として扱い，
命題\ref{ke:collapsed-tree-common-moment}の共通 moment parameter 化と合わせて，
命題\ref{ke:collapsed-tree-full-fiber-collapse}で潰される木上の全 fiber 合併径を得る．
任意に異なる moments に属する neck 全体の径が消えるという主張ではない．
全曲面の no-shortcut と距離上界は，定理\ref{ke:global-ambient-lower}，
\ref{ke:global-ambient-upper}で別に証明する．

\begin{corollary}[Central source を含む深部での短い円]
\label{ke:source-short-circle}
Source の固定正則座標 $xy=tz$ で，十分小さい固定 $r>0$ を取る．
$\rho_a(x,y,z)=(e^{2\pi\sqrt{-1}a}x,e^{-2\pi\sqrt{-1}a}y,z)$ の
周期 $1$ の生成場を $V_{\rm src}$ とする．
$|x|,|y|,|z|\le r$ の範囲で，$L_t=-\log|t|^2$ として
\[
 |V_{\rm src}|_{g_t}
 \le \frac{C}{\min\{-\log|x|^2,-\log|y|^2,L_t\}},
\]
ただし零の座標の深度は $+\infty$ とする．特に
$|x|,|y|\le c|t|^{1/2}$ の central source で短さは $O(L_t^{-1})$ である．
\end{corollary}

\begin{proof}
小さい固定 $c>0$ について，$x$-oriented 部分
$|x|\ge c|t|^{1/2}$ を考える．$(x,z)$ を用いる annulus--disc extension
$y=(t/x)z$ に対し，関連 KE 主稿
\texttt{coarse:lem:pair-arm-density} の capped density 比較と
\texttt{coarse:lem:complete-volume-upper} により $\varphi_t\le C$ を得る．
実際，$r$ は extension の人工境界から固定幅を取って選び，
$s_x=-\log|x|^2+O(1)\le L_t/2+O(1)$ なので，
$\alpha_y\asymp L_t$，$\alpha_x\asymp s_x$ であり，
capped density と product envelope は比較可能である．
これには gate $|x|=c|t|^{1/2}$ 上も含まれる．
補題\ref{ke:capped-projective-lower}を使うと
$g_t\ge c' g_{\mathrm{FS}}$ が成立する．固定した ambient coordinate $z$ の
微分は射影計量に対して有界なので，残余 rank-one 下界
$p_z^*h_D\le C\omega_t$ が成り立つ．
体積上界と \texttt{neck:lem:rank-one-circle} を使い，
$x$ の角速度を持つ $V_{\rm src}$ の長さを $C/s_x$ 以下に抑える．
$y$-oriented 部分も対称である．

残る領域は $(x,y)$ で表した rescaled bidisc
$\{|x|,|y|<c|t|^{1/2}\}$ である．$c^2<r$ とすれば
その閉包は $|z|\le c^2$ で source chart 内に含まれる．
境界では少なくとも一方の oriented 評価により
$|V_{\rm src}|_{g_t}=O(L_t^{-1})$ である．
$V_{\rm src}$ は正則ベクトル場
$W=x\partial_x-y\partial_y$ の虚方向の実生成場であり，
$|V_{\rm src}|_{g_t}^2$ は $|W|_{g_t}^2$ の固定倍である．
正則ベクトル場の Bochner 恒等式
$\Delta_{\omega_t}|W|^2=|\nabla^{1,0}W|^2-\operatorname{Ric}(W,\bar W)\ge0$
と最大値原理で，同じ上界が bidisc 内部にも成り立つ．
\end{proof}

この系は source の円作用そのものを評価し，固定点での作用の退化を許す．
移動する normal 境界までの actual buffer と閉じた pinched torus の構成は，
次の系\ref{ke:source-diagonal-moser}および定理\ref{ke:projected-source-neck}で扱う．
Source 中央帯の全点径は定理\ref{ke:source-central-diameter}，
全共通 interval の同一 moment 退出は定理\ref{ke:source-all-moment-exit}で別に評価する．

\begin{lemma}[Source の零 moment 商と閉じた pinched fibre]
\label{ke:source-reduction-pinched-fiber}
$t\ne0$ を固定し，境界付き compact 残余曲面 $P$ の近傍に
\[
 X=\{xy=t h(p)\},\qquad
 \rho_a(x,y,p)=(e^{2\pi\sqrt{-1}a}x,e^{-2\pi\sqrt{-1}a}y,p)
\]
を取る．$h$ の零点は $P$ の内点にある単純零点とし，
$\omega=d\lambda$ は $X$ 上で滑らかな K\"ahler form とする．
$\eta=\mathcal A\omega$，$\bar\lambda=\mathcal A\lambda$，
$\mu=\bar\lambda(V)$ と置く．固定 $R>0$ と
$0\in\operatorname{int}I$ なる閉区間 $I$ を選び，
$|t|\sup_P|h|<R^2$，および全 $p\in P$ で
\[
 \left.\mu\right|_{|y|=R}<\inf I<\sup I
       <\left.\mu\right|_{|x|=R}
\]
を仮定する．$K=\{p\in P,\ |x|,|y|\le R,\ \mu\in I\}$ の
飽和開近傍で centered normal Moser flow が全時刻存在するとする．
このとき次が成り立つ．
\begin{enumerate}
 \item 各 source $o=(0,0,q)$ で $\mu(o)=0$ であり，
 centered Moser flow は $o$ を固定する滑らかな微分同相である．
 Actual circle の symplectic isotropy representation は weights $(1,-1)$ を保ち，actual moment は
 $\widehat\mu=\mu\circ\phi^{-1}$ である．その値は actual 軌道の
 $\lambda$-周期に等しく，source ごとに加法定数を選び直す必要はない．
 \item $(\mu,p)$ は $K/S^1\cong I\times P$ を与える．
 $(0,q)$ 以外ではこれは滑らかな orbit space の座標で，
 $K\setminus\{o:h(q)=0\}$ 上の作用は自由である．
 \item $P$ の marked Morse foliation で，各 source が
 saddle を含まない相異なる閉 regular leaves 上にあるとする．
 $\mu$ を固定して葉を持ち上げた閉 fibre は，通常 $T^2$，
 単純 saddle leaf では $S^1\times\figeight$，
 $\mu=0$ で source を一つ含む葉では一重 pinched torus となる．
 二次元 smooth strata は Lagrangian で，各 source は
 nondegenerate focus--focus point である．近傍の regular torus 上の
 normal circle は primitive な vanishing cycle を与える．
 これらの結論は $\phi$ によって actual $\omega$ へ移される．
\end{enumerate}
ここでの滑らかさは各固定 $t\ne0$ に対する主張であり，
$t\to0$ での flow の導関数や正規形の一様性は含めない．
\end{lemma}
\begin{proof}
Source の近くでは $h'(q)\ne0$ により $(x,y)$ が $X$ の正則座標である．
平均 $\eta,\bar\lambda$ と
$\beta=\int_0^1(1/2-a)\rho_a^*(\iota_V\omega)\,da$ はここでも滑らかである．
$V(o)=0$ なので $\mu(o)=\beta(o)=0$．
$\omega_s=(1-s)\eta+s\omega$ は全時刻で K\"ahler form であり，
$\iota_{Y_s}\omega_s=\beta$ は滑らかな場 $Y_s$ を与え，$Y_s(o)=0$ となる．
補題\ref{ke:circle-exact-transport}の証明は自由性を使わないので，
固定点を含めてその exact identity と周期の結論が成立する．
共役によって固定点での円表現も変わらない．

$h(p)\ne0$ の fibre では，$x\mapsto e^r x$，$y\mapsto e^{-r}y$ の
実動径場は $-JV/(2\pi)$ である．従って
\[
 d\mu\bigl(-JV/(2\pi)\bigr)=|V|_\eta^2/(2\pi)>0.
\]
各 moment level はただ一つの円となる．$h(p)=0$ では二軸が原点を共有し，
$x$ 軸で $\mu$ は零から正へ，$y$ 軸で零から負へ狭義単調に変化する．
Cut の仮定により各 $(\mu,p)\in I\times P$ の逆像は一軌道である．
$K/S^1$ の compact 性から連続全単射は同相となる．
固定点の外では $dp$ は complex rank $1$ を持ち，
$dp(JV)=0$，$d\mu(JV)=-|V|_\eta^2\ne0$ であるから，
$(\mu,p)$ は滑らかな submersion である．Weights が $(1,-1)$ なので
固定点以外に有限 stabilizer もない．

持ち上げた葉の smooth 二次元接空間は $V$ と
$\ker d\mu$ 内の一方向で張られ，$\iota_V\eta=-d\mu$ より isotropic である．
固定点を含まない円及び saddle graph 上の主円束は，
基底の $H^2$ が零なので自明であり，所記の fibre 型を得る．
Source $q$ を一つ含む円葉 $\gamma$ では，
$\gamma\setminus\{q\}$ 上の主円束は $S^1\times(0,1)$ である．
閉 fibre は compact Hausdorff で，その補集合はただ一つの固定点である．
従ってこれは cylinder の一点 compactification，すなわち
$S^2$ の二点を同一視した一重 pinched torus である．

残余葉の局所 submersion を $a(p)$ とする．
これは円不変なので $\mu$ と Poisson 可換である．
Source で $\eta(o)$ は正の Hermitian 形式であり，円不変性から
$(x,y)$ の二つの異なる weights に関して対角である．
複素線形 rescaling 後，moment の二次項は
$\pi(|u|^2-|v|^2)$ となる．
$da(q)\ne0$ と $p-q=xy/(t h'(q))+O((xy)^2)$ により，
$a(p)-a(q)$ の二次項は $\Re(cuv)$，$c\ne0$ である．
Phase 変更と実線形な基底値の再規格化により，この組を
$F=(|u|^2-|v|^2)/2$，$H=\Re(uv)$ に直す．
$u=\alpha+\sqrt{-1}\beta$，$v=\gamma+\sqrt{-1}\delta$ と書くと，
\[
 (Q_1,P_1,Q_2,P_2)
 =2^{-1/2}(\alpha+\delta,\,\beta-\gamma,\,\beta+\gamma,\,\delta-\alpha)
\]
は $d\alpha\wedge d\beta+d\gamma\wedge d\delta=\sum_{j=1}^2dQ_j\wedge dP_j$ を満たし，
\[
 F=-(Q_1P_2-Q_2P_1),\qquad H=-(Q_1P_1+Q_2P_2)
\]
となる．従ってこれは focus--focus の標準二次型
\cite[Section~2]{ZungFocus}であり，局所的な node は非退化で，
零次元 stratum は固定点そのものである．
Regular torus は円上の主円束なので，normal 軌道はその
$H_1(T^2;\mathbb Z)$ の primitive class を表す．
$p=q$ 上で $\mu\to0$ とすると軸上のこの円が固定点へ縮むため，
それが vanishing cycle である．重複被覆による multiple cycle は生じない．
最後に $\phi^*\omega=\eta$ を用いて全ての strata と circle を移す．
\end{proof}

\begin{remark}[零 reduction の標準模型と ambient の滑らかさ]
標準 product
$\omega_0=(\sqrt{-1}/2)(du\wedge d\bar u+dv\wedge d\bar v)$，
$\mu=\pi(|u|^2-|v|^2)$，$w=uv$ では，$w\ne0$ 上の reduced form は
\[
 \omega_{\mathrm{red},\mu}
 =\frac{\sqrt{-1}\,dw\wedge d\bar w}
        {2\sqrt{(\mu/\pi)^2+4|w|^2}}.
\]
実際，$m=\mu/(2\pi)$，$\rho=|w|$ とすれば
$(|u|^2+|v|^2)/2=\sqrt{m^2+\rho^2}$ であり，
reduction は $\frac12 d\sqrt{m^2+\rho^2}\wedge d\arg w$ となる．
$\mu=0$ の商は位相的には $w$-平面だが，その symplectic density は
$w=0$ で $|w|^{-1}$ の特異性を持つ．この商を通常の滑らかな
symplectic 曲面とみなす必要はなく，前補題は ambient の
滑らかな形式と接空間を用いて Lagrangian 性を証明している．
この明式は標準模型の計算であり，任意の平均された $\eta$ に対して
Darboux 座標が元の residual 座標 $p$ 又は $w=t h(p)$ を保つとは主張しない．
Focus--focus の局所型と monodromy は既知の理論であり，ここで行ったことは
固定点を含む actual KE neck の exact transport と閉 fibre への適用である．
\end{remark}

\begin{corollary}[Moment に依存する残余関数の focus 判定]
\label{ke:source-focus-transversality}
前補題の固定 source $o=(0,0,q)$ で，$A(\mu,p)$ を通常の実座標
$(\mu,\Re p,\Im p)$ の滑らかな関数とする．
可換な写像 $F=(\mu,A(\mu,p))$ 自身が $o$ で非退化 focus--focus 型であるための
必要十分条件は $d_pA(0,q)\ne0$ である．
\end{corollary}
\begin{proof}
前補題の正規化を用いると，二次項は
\[
 Q_\mu=\pi(|u|^2-|v|^2),\qquad
 Q_A=A_\mu(0,q)Q_\mu+\Re(cuv)
\]
であり，$c\ne0$ と $d_pA(0,q)\ne0$ は同値である．
このとき $Q_\mu$ の倍を引いて前補題の focus 二次型を得る．
$c=0$ なら二次形式の span は高々一次元なので非退化ではない．
\end{proof}

この系は指定した momentum map の判定である．例えば残余 submersion $a$ に対し
$A=(a-a(q))^3$ と置くと同じ葉の分割を与えるが，写像の二次項は退化する．
本稿の source 近傍では $a$ 自身を用い，その germ を後続の相対補間で固定する．

\begin{corollary}[Source を含む遅い exhaustion 上の actual 円作用]
\label{ke:source-diagonal-moser}
上の source chart で固定した十分小さい $b>0$ を取る．
$r_t\to0$ を十分遅く選べば，$xy=tz$ の飽和集合
$K_{t,r_t}=\{|x|,|y|\le r_t^2,\ |z|\le b/2\}$ は
normal Moser map の source と像の双方に含まれる．変位と円軌道像の
ambient 直径は $O(1/|\log r_t|)$ で，source の固定点は移らない．
この近傍に actual $\omega_t$ に対する focus--focus germ が得られる．
\end{corollary}

\begin{proof}
まず $r$ を固定して $U_{t,r}=\{|x|,|y|<r,\ |z|<b\}$ とする．
$K_{t,r}$ からの最初の退出が $|z|=b$ で起こる場合，経路は
例えば $3b/4\le |z|\le b$ の固定 shell を横切る．これは source を避ける
固定残余 compact に属し，十分小さい $r$ と $t$ で
定理\ref{ke:fixed-base-neck-moser}の証明中の射影下界（又は局所積比較）が
$z$ の変位に一定の actual 費用を与える．この shell では元の $z$ 自体が
source-free pair の残余座標となる．

$|x|=r$ で退出する場合は $r^2\le|x|\le r$ を横切る．
補題\ref{ke:source-inclusive-fixed-shell}の shell 評価を，固定 source 座標
$xy=tz$，$|z|\le b$ に適用する．各 $r$ を先に固定してから全ての
十分小さい複素 $t\ne0$ を取れば，この横断費用は $r$ に独立に正である．
$z=0$ も同補題の固定成分 compact に含まれる．$y$ も同様である．

よって各固定 $r$ について $|t|$ を十分小さくすれば，$K_{t,r}$ と
$U_{t,r}$ の境界の間の actual buffer は $r$ に独立な正の定数以上である．
補題\ref{ke:short-circle-actual-buffer}と系\ref{ke:short-circle-image-buffer}は
作用の自由性を使わないので，$U_{t,r}$ 上の
$|V_{\rm src}|=O(1/|\log r|)$ に適用できる．
必要な比較的 compact な中間領域は，境界を少し内側に動かすか，
$K_{t,r}$ の十分小さい actual 距離近傍を用いて取れる．
$r_n\downarrow0$ と，各 $r_n$ に対する閾値 $t_n\downarrow0$ を選び，
$t_{n+1}<|t|\le t_n$ では $r_t=r_n$ とする．
これで被覆と変位上界を得る．$(x,y)$ は $X_t$ 上の正則座標であり，
$dx\wedge dy$ に対する対数体積は滑らかな $dd^c$-primitive を与える．
補題\ref{ke:source-reduction-pinched-fiber}の固定点での transport と
局所二次型の議論を $h(p)=p=z$，残余 submersion $\Re z$ に適用すれば，
固定点の保存と actual focus--focus germ を得る．
\end{proof}

この diagonal 選択は所定の $t$ 依存の速さを保証しない．
外側の foliation と接続する閉じた pinched torus は
定理\ref{ke:projected-source-neck}で構成し，
同一 moment の source-free 境界への接続費用は定理\ref{ke:source-all-moment-exit}で与える．

\subsection{Source を含む一つの射影的 neck}

\begin{lemma}[射影を保つ source-inclusive normal form]
\label{ke:projected-source-normal-form}
固定二重線 $D=D_{ij}$ 上で，三重点を除く固定 compact 領域の近傍を取る．
関連 KE 主稿 \texttt{pair:prop:bundle-conic} と同じ bundle projection を用い，
$h=-F_d/\prod_{k\ne i,j}L_k$ とする．固定幅の零切断近傍で
射影を保つ正則座標変換があり，退化族の式は
\[
 x'y'=t h_t(p),\qquad h_t=h|_D+O(t)
\]
となる．ここで $x',y'$ は元の二つの直線束に値を持ち，座標変換は
恒等写像から全固定階で $O(t)$ ずれる．$h_t$ の零点は旧 source の近傍に
一つずつある単純零点である．従って全基底上で一つの正則円作用
$(x',y')\mapsto(e^{2\pi\sqrt{-1}a}x',e^{-2\pi\sqrt{-1}a}y')$ を得る．

より正確には，基底と Hermitian norms を先に固定し，
$E_a=\{\|x\|<a,\ \|y\|<a\}$ と置くと，正数 $r,\epsilon$ を選べる．
座標変換 $\Phi_t$ は $|t|<\epsilon$ を含めて正則に依存し，
$E_{2r}$ 上で定義され，その像は $E_r$ を含む．$E_r$ 上の逆写像と共に，
各固定階 $q$ に対して
\begin{equation}\label{ke:eq:source-normal-form-common-domain}
 \|\Phi_t-\id\|_{C^q(E_{2r})}
 +\|\Phi_t^{-1}-\id\|_{C^q(E_r)}\le C_q|t|
\end{equation}
を満たす．Norm は固定した有限個の束 charts で測り，必要なら外側の
tube を最初に少し縮める．円作用の定義域は
$\Phi_t^{-1}(E_r)\cap X_t$ と取れば，その全 orbit が同じ定義域に入る．
各 source は $t=0$ を含む円板上の一価な正則断面として追跡でき，
$t\ne0$ ではこれらの断面が円作用の固定点集合である．ここでは基底を三重点へ動かした
ときの $r,\epsilon,C_q$ の一様性は主張しない．
\end{lemma}

\begin{proof}
束の値を持つ関数 $f_t(p,x,y)=xy-th(p,x,y)$ の fibre derivative は
$t=0$ の零切断で非退化な Hessian を持つ．陰関数定理から一意な
正則臨界断面 $a_t=O(t)$ を得る．臨界点の一意性により線形 frame の間で貼り合う．
$\xi=(x,y)-a_t$ とし，
$f_t-f_t(p,a_t)=\xi^{\mathsf T} Q_t(p,\xi)\xi$ と Taylor 展開する．
$Q_t$ は Hessian の線分積分で取れる対称な行列であり，
$Q_t=Q_0+O(t)$，$\xi^{\mathsf T}Q_0\xi=\xi_1\xi_2$ である．
$M=Q_0^{-1}Q_t$ は $Q_0$ に関し自己共役である．恒等行列の近傍で
二項級数により $S=M^{1/2}$ を取ると $S^{\mathsf T}Q_0 S=Q_t$．
Frame 変更で $M,S$ は共役変換するので，$\xi'=S\xi$ は束全体で定義される．
固定幅を縮めれば $\xi\mapsto\xi'$ は一様に恒等写像に近い微分を持つ正則同型である．
$t h_t=-f_t(p,a_t)$ とすれば所要の式を得る．
$a_t=O(t)$ により $h_t=h|_D+O(t)$ となり，零点は陰関数定理と
零点近傍の外の一様な非消滅から従う．

量化と定義域を補足する．基底 compact を含む少し広い固定近傍と有限個の
束 charts を先に選ぶ．その上で Hessian の非退化性，臨界断面の陰関数定理，
および $M$ の恒等行列からの距離を一様に取れる．平方根は
$\|M-I\|<1/2$ における同じ収束二項級数で定義するので，
$t$ の argument による枝の変更はない．各 frame で
$\Phi_t(z)=z+B_t(p,z)$ と書けば，外側の固定 tube 上で
$B_t=O(t)$ が全固定階で成立する．$|t|$ を小さくした後は
$\|D_zB_t\|<1/2$ であり，内側 tube の任意の $w$ に対する
$z=w-B_t(p,z)$ は外側 tube 内の固定小球で一意に解ける．
これが共通像と逆写像を与え，微分した同じ方程式から
\eqref{ke:eq:source-normal-form-common-domain}を得る．
Frame での一意な逆写像は重なりで一致する．

$f_t(p,a_t)$ は $t=0$ で恒等的に零なので，$h_t$ は $t=0$ にも正則に延びる．
各単純零点 $q$ に対する $h_t(p)=0$ の陰関数定理は，
一つの円板 $|t|<\epsilon$ 上で $q(t)=q+O(t)$ を与える．
円板を共通に縮めればこれらは互いに交わらず，他に零点はない．
標準回転の固定点は $x'=y'=0$ であり，$t\ne0$ では族の式から $p=q(t)$ に限られる．
元の座標では $(p,x,y)=(q(t),a_t(q(t)))$ である．
最後に $E_r$ は二つの束 norms により定義したため標準回転で不変であり，
$\Phi_t^{-1}$ による移送は定義域全体での円作用を与える．
従って部分 orbit から平均を定義する必要はない．
\end{proof}

\begin{lemma}[Source を横切る固定 end と shell]
\label{ke:source-inclusive-fixed-shell}
固定次数 $d\ge5$ の本稿の一般位置全平面退化を考える．三重点を避けた
二重線上の固定 compact $P\Subset\operatorname{int}P^+$ を取り，
補題\ref{ke:projected-source-normal-form}の一つの affine frame で
$x'y'=t h_t(p)$ と書く．$P^+$ は source を含んでよい．固定 source 座標
$xy=tz$ における同じ結論も含める．$V$ は $x'$ を正方向，$y'$ を負方向に
回す周期 $1$ の生成場，$\mathcal A$ はその円平均とする．
次の一つの体積形式と primitive を用いる：
\begin{equation}\label{ke:eq:source-inclusive-fixed-frame}
\begin{aligned}
 \Omega_t=d\log x'\wedge dp=-d\log y'\wedge dp,
 \\
 \lambda_t=d^c\log\frac{\omega_t^2}{|\Omega_t|^2},
 \qquad \mu_t=(\mathcal A\lambda_t)(V).
\end{aligned}
\end{equation}
最初の二つの表示はそれぞれ $x'\ne0$，$y'\ne0$ 上で読む．
この $\Omega_t$ は source にも非零正則に延び，$d\lambda_t=\omega_t$ である．

ある $\delta>0$ が存在し，十分小さい各固定 $r>0$ と $s_r=-\log r^2$ に対し
\begin{align}
 \mu_t|_{\{p\in P,\ |x'|=r\}}
   &=\frac{4\pi}{s_r}+O_{P^+}(s_r^{-1-\delta})+o_{t\to0;r,P^+}(1),
       \label{ke:eq:source-inclusive-positive-end}\\
 \mu_t|_{\{p\in P,\ |y'|=r\}}
   &=-\frac{4\pi}{s_r}+O_{P^+}(s_r^{-1-\delta})+o_{t\to0;r,P^+}(1)
       \label{ke:eq:source-inclusive-negative-end}
\end{align}
が全 base points と全 phase に一様に成り立つ．$O_{P^+}$ の定数は小さい $r$ に
共通であり，収束誤差を抑える $t$ の閾値は $r$ に依存してよい．
さらに $c_{P^+}>0$ を $r$ に独立に選べる．各固定 $r$ に対して十分小さい全 $t$ では，
\[
 \Sigma^x_{t,r}=\{p\in P^+,\ r^2\le|x'|\le r,\ y'=t h_t(p)/x'\}
\]
内で $|x'|=r^2$ と $|x'|=r$ を結ぶ全経路の actual 長さは
$c_{P^+}\log2$ 以上である．$x',y'$ を交換した shell でも同じ結論が成り立つ．
\end{lemma}

\begin{proof}
まず固定成分への収束の適用条件を確認する．全空間は平坦な射影 hypersurface family で，
relative canonical bundle は adjunction により
$K_{\mathcal X/\Delta}=\mathcal O_{\mathcal X}(d-4)$ であり，Cartier かつ relatively ample である．
中央 fibre は一般位置平面の reduced normal-crossing union なので semi-log canonical である．
全空間の特異点は source の $xy=tz$ のみで，三次元 ordinary double points である．
原点の blowup の discrepancy は $3-2=1$ なので全空間は normal かつ canonical である．
従って Song~\cite[Definition~2.3, Lemmas~4.4--4.5]{SongNegativeKE}の stable degeneration
に対する固定 compact の $C^\infty$ 評価と極限の一意性を適用できる．
関連 KE 主稿\cite{KELocalCurrent}の \texttt{pair:thm:fixed-end-moment} の証明と同じく，
各 component 上の極限は $K_{\PP^2}+D=\mathcal O_{\PP^2}(d-4)>0$ に対応する
complete cusp KE 計量である．この同定には Song~\cite[Lemma~3.7]{SongNegativeKE}の
一意性と，cusp potential が任意の正係数の logarithmic pole より緩やかなことを用いる．

$r$ を固定すると
\[
 (p,x')\longmapsto(p,x',t h_t(p)/x')
\]
は $P^+\times\{r^2\le|x'|\le r\}$ の少し広い近傍を，$t=0$ まで滑らかに
同定する．極限の $x'$ は零でないから，その像は中央成分 $y'=0$ の smooth locus の
固定 compact に入る．特に $h_0(p)=0$ でも中央 fibre の二重 locus を踏まない．
射影 Morse 変換は全固定階で $O(t)$ なのでこの固定 gauge の収束を保つ．
Song の compact 評価から任意の $t_\nu\to0$ に収束部分列があり，極限の一意性から
全複素 parameter $t\to0$ で同じ $C^\infty$ 収束を得る．
従ってこの領域上の $\log(\omega_t^2/|\Omega_t|^2)$ とその全固定階微分も，
$d\log x'\wedge dp$ で測った component の対数体積 $F_0$ へ収束する．

この成分上の唯一の局所 boundary divisor は $x'=0$ であり，source を指定する
$h_0(p)=0$ は追加 divisor ではない．$q=x'$，$s=-\log|q|^2$ とすると，
Rochon--Zhang~\cite[\S5, (5.5)--(5.6), Theorem~5.1]{RochonZhang}の表示は
\[
 F_0=-2\log(s+b(p,q))+B(p,q)+\upsilon(p,q),\qquad
 \upsilon-\Xi\upsilon_D\in s^{-\nu}C^\infty_{\rm fc},\quad \nu>0
\]
となる．$b,B$ は $q=0$ を跨いで滑らかである．固定 $P^+$ は他の boundary divisors
から離れているので，$\Xi\upsilon_D$ の radial derivative は
$O_{P^+}(s^{-2})+O_{P^+}(|q|)$ である．また $|q|\partial_{|q|}$ の cusp norm は
$O(s^{-1})$ だから，remainder の一階評価を合わせて
\[
 |q|\partial_{|q|}F_0=\frac4s+O_{P^+}(s^{-1-\delta}),\qquad
 \delta=\min\{\nu,1\}>0
\]
を得る．有限個の fixed patches と二つの components で定数と指数を共通に取る．
周期 $1$ の正の回転 $V_q$ について
$(d^cF_0)(V_q)=\pi |q|\partial_{|q|}F_0$ であり，円平均もこの一様評価を保つ．
$x'$-end では $V=V_{x'}$，$y'$-end では $V=-V_{y'}$ なので，固定 gauge の
収束誤差を加えて両 end の式を得る．

同じ component cusp 比較から，固定 $P^+$ 上の十分小さい全 $|x'|>0$ で
$g_0\ge c\,(d\log s)^2$ となる $c>0$ を取れる．固定 shell 上での smooth convergence により，
各 $r$ に対して $t$ を十分小さくすれば $g_t\ge(c/2)(d\log s)^2$ である．
Shell の両端では $s=2s_r,s_r$ なので，任意の横断経路の長さは
$\sqrt{c/2}\,|\log(2s_r)-\log s_r|=\sqrt{c/2}\log2$ 以上である．
これは shell 内の部分経路の評価であり，任意の全曲面経路が shell を横切る場合にも
その部分へ適用できる．Shell 外を許す二点間の距離を直接評価する主張ではない．

最後に primitive の規格化を確認する．$x'y'=t h_t(p)$ の微分を $dp$ と wedge すると，
二つの logarithmic 表示が重なりで厳密に一致する．Source 近傍では単純零点の仮定より
\[
 \Omega_t=\frac{dx'\wedge dy'}{t h'_t(p)}
\]
と非零正則に延びる．よって一つの $\lambda_t$ が全域に定義され，負 KE 方程式から
$d\lambda_t=\omega_t$ を得る．Source に追加の pole や留数補正は生じない．
この主張は今固定した同一の $\Omega_t$ に関するものである．別の非零正則 frame
$\widehat\Omega_t=b\Omega_t$ に変えれば primitive の差は
$-d^c\log|b|^2$ という閉形式となる．特に $b=b(p)$ なら normal moment への寄与は零だが，
residual periods が零とは限らない．異なる正規形の座標写像から定めた primitive の
無条件な一致も主張しない．
\end{proof}

\begin{theorem}[固定基底上の共通 actual neck と閉じた特異ファイバー]
\label{ke:projected-source-neck}
$P$ を三重点を避けた固定 compact な境界付き残余曲面とし，
source は内点にあり，境界近傍にはないとする．
$P$ 上に閉じた regular circle leaves と単純 saddle leaves を持つ
marked Morse foliation を与える．各 saddle leaf は一つの saddle を含み，
各 source は saddle のない相異なる regular leaves に置くとする
（動く source に合わせた小さい isotopy を許す）．
このとき，十分小さい $|t|$ について，ある共通 moment interval $I_t$ 上で
同じ円作用を用いる actual $\omega_t$-Lagrangian neck fibration が存在する．
Regular fibre は $T^2$，saddle fibre は $S^1\times\figeight$，
各 source の critical fibre は一重 pinched torus である．
\end{theorem}

\begin{proof}
$P$ を少し広げた固定基底 $P^+$ に補題\ref{ke:projected-source-normal-form}を適用する．
三重点の一つを無限遠とする affine 座標 $p$ を取り，二つの直線束を
自明化する．以下の絶対値はこの固定 frame による．
$r>0$ を固定し，外側を $p\in\operatorname{int}P^+$，$|x'|,|y'|<r$，
内側を $p\in P$，$|x'|,|y'|\le r^2$ とする．
Source の近傍では $z=h_t(p)$ が正則残余座標で
$x'y'=tz$ となる．この $t$ 依存座標でも系\ref{ke:source-short-circle}の
oriented arm の体積比較は一様である．実際，oriented 側
$|x'|\ge c|t|^{1/2}$ では元の平面座標 $x=x'+O(t)$ と $x'$ が比較可能で，
$s_x=-\log|x'|^2+O(1)$ である．もう一方は $y'=tz/x'$ と $y=y'+O(t)$
から $|y|=O(|t|^{1/2})$ である．Canonical residue は
$d\log x'\wedge dz$ の一様に有界な非零倍である．これは
$f_t$ と元の方程式の比が非零 unit で，座標変換の Jacobian と
$h'_t(p)$ が一様に非退化であることによる．従って capped density は
arm の annulus--disc envelope と同じ率を持つ．反対側も同様で，
central bidisc には Bochner 最大値原理を同じ作用に適用する．
Source を避ける固定部分では $h_t$ は一様に非消滅であり，
$x'y'=th_t(p)$ は $y'/h_t(p)$ によって固定 pair の積座標になる．
各 cap は元の座標と新座標の絶対値を $|t|^{1/2}$ で切ったものに関して
一様比較可能なので，ordinary pair の体積比較と
定理\ref{ke:fixed-base-neck-moser}の証明が適用できる．
有限被覆で外側領域全体の生成場長は $O(1/|\log r|)$ となる．

残余方向の最初の退出には $P^+\setminus P$ 内の source-free 固定幅を
横切る必要があり，そこでの射影下界によって一定の費用がある．
法線方向の退出については $r^2\le|x'|\le r$ 又は
$r^2\le|y'|\le r$ の shell を横切る．それぞれ $r$ を固定して
$t\to0$ とすれば，$P^+$ 全体で成分の smooth locus に固定 gauge で
収束する．座標変換は $O(t)$ であり，補題\ref{ke:source-inclusive-fixed-shell}の
smooth convergence と cusp 下界により，この横断費用も $r$ に一様に正である．
$ r=r_t\to0$ を対角的に選び，系\ref{ke:short-circle-image-buffer}を
適用して内側を source と像の両方に含む一つの Moser map を得る．
固定点を許すことは系\ref{ke:source-diagonal-moser}と同様である．

正則体積 frame $\Omega=d\log x'\wedge dp$ は $y'\ne0$ で
$-d\log y'\wedge dp$ と一致し，source では
$dx'\wedge dy'/(t h'_t(p))$ として非零に延びる．
従って $\lambda_t=d^c\log(\omega_t^2/|\Omega|^2)$ を全体で用いられる．
$\eta=\mathcal A\omega_t$ とし，$\mu=(\mathcal A\lambda_t)(V)$ を取る．
ここで normal Moser map には centered kernel $1/2-s$ を用いる．
同じ選択を三重点側にも用い，補題\ref{ke:normal-flow-naturality}で
共通の作用に関する flow の自然性を適用する．
内側 cutoff の $|x'|=r_t^2$ と $|y'|=r_t^2$ における $\mu$ は逆符号で，
絶対値は $c/|\log r_t|$ 以上ある．補題\ref{ke:source-inclusive-fixed-shell}に
半径 $r^2$ を代入し，まず各 $r$ で固定成分の remainder を主項より小さくしてから，
全複素 $t$ に対する収束誤差も同時に小さくする対角選択を用いる．
これには $h_t=0$ 上の cutoffs と全 phase が含まれ，primitive は上記の
一つの $\Omega$ から定めたもののままである．
さらに対角閾値を縮めて $|t|\sup_{P^+}|h_t|<r_t^4$ とし，
$I_t$ は両 cutoff の値の間にある $0$ を含む閉区間に取る．

これで補題\ref{ke:source-reduction-pinched-fiber}の仮定が
$R=r_t^2$，$h=h_t$，$I=I_t$ に対して満たされる．
同補題の商 $I_t\times P$ 上で残余 foliation の葉を取り，
exact Moser map で移せば，閉 fibre 全体の Lagrangian strata，
連結性と三種類の位相型を得る．各 source の vanishing cycle は
共通 normal circle の primitive class であり，その critical moment は正確に零である．
\end{proof}

ここで構成した neck は三重点を除く固定残余基底に沿う．
三重点 collar における同一の fiber partition は
系\ref{ke:neck-corner-saturated-union}，
特異 fibers を含む全 fiber 合併の正規化 ambient 直径評価は
命題\ref{ke:collapsed-tree-full-fiber-collapse}で別に与える．

\subsection{閉じた focus fibre の period と近傍全体の移送}

Focus 点の局所正規形と，compact singular fibre の飽和近傍全体の
symplectic germ を区別する．後者の Taylor 不変量と，primitive を指定した場合の
追加の定数項は V\~u Ng\d{o}c\cite[\S\S3, 7]{VuNgoc}の既知の不変量である．
以下では本稿の exact transport に必要な period 条件を直接記す．

\begin{lemma}[一重 focus 近傍での一つの残余 period]
\label{ke:focus-neighborhood-period}
$F:N\to D$ を円板上の proper な Lagrangian torus fibration とし，
唯一の特異 fibre $\Lambda=F^{-1}(0)$ は一つの非退化 focus 点 $o$ を持つ
compact connected pinched torus とする．$D$ を十分小さく取れば
\[
 H_1(N;\ZZ)\cong\ZZ,\qquad H_2(N;\ZZ)\cong\ZZ.
\]
また $\omega|_N$ は exact である．
正則 torus $T$ 上で，primitive な vanishing cycle を $\alpha$，
それを含む整数基底を $(\alpha,\beta)$ とすると，包含写像は
\[
 H_1(T;\ZZ)/\ZZ\alpha\ \xrightarrow{\ \sim\ }\ H_1(N;\ZZ)
\]
を誘導する．従って $N$ 上の滑らかな closed one-form $\zeta$ に対し
\begin{equation}\label{ke:eq:focus-residual-exactness}
 \zeta\text{ が exact}\quad\Longleftrightarrow\quad
 \int_\beta\zeta=0.
\end{equation}
Normal period $\int_\alpha\zeta$ は自動的に零である．
特に一つの外側の正則 fibre $T=F^{-1}(c)$，$c\ne0$ の近傍で $\zeta=0$ なら，
$\zeta$ は $N$ 全体で exact である．
\end{lemma}
\begin{proof}
Eliasson の局所正規形に適合する複素表示で
$q_1+\sqrt{-1}q_2=\bar z_1z_2$ とする
\cite[\S2]{ZungFocus}．局所の base 座標変更を用い，
$|z_1|,|z_2|<r$，$|\bar z_1z_2|<\delta\ll r^2$ で定める node block $B$ を取る．
同時の radial contraction が $B$ を原点へ縮めるので $B$ は可縮である．
二つの boundary faces はそれぞれ円板上の circle bundle であり，
corner $|z_1|=|z_2|=r$ は $|q|<\delta$ に現れない．
これらに沿って少し重ねた外側の regular block を $C$ とする．
$\Lambda$ の $B$ 外の部分は compact annulus である．
Base をさらに縮め，閉じた annulus block 上の boundary を含む proper submersion に
Ehresmann の議論を適用する．その collar を少し厚くした開集合を $C$ とすれば，
円板上の自明な annulus bundle と同じ homotopy type を持つ．従って
\[
 B\simeq\{*\},\qquad C\simeq S^1,
 \qquad B\cap C\simeq S^1\sqcup S^1.
\]
Mayer--Vietoris の写像 $H_1(B\cap C)\to H_1(C)$ は，
向きを揃えれば $(m,n)\mapsto m-n$ という全射である．
その kernel と，$H_0(B\cap C)\to H_0(B)\oplus H_0(C)$ の kernel はともに $\ZZ$ なので，
表示した $H_2$ と $H_1$ を得る．
$H_2$ の generator は，$B$ 内の二つの局所 Lagrangian discs と
$C$ 内の annulus からなる pinched fibre の二次 cycle である．
その上の $\omega$ の積分は零なので，de Rham の定理より $\omega|_N$ は exact である．
Normal circle は $B$ 内で潰れ，$\beta$ は二つの overlap components を
$B$ と $C$ を通って一周するので後者の kernel の generator へ写る．
これで包含写像の主張と，de Rham の定理による\eqref{ke:eq:focus-residual-exactness}が従う．
外側の任意の regular torus 上でも同じ $\beta$ が generator である．
\end{proof}

\begin{example}[Node の primitive germ を固定しても残る定数]
\label{ke:focus-local-primitive-shift}
前補題の regular block を $C\cong S^1\times(0,1)\times D$ と書き，
$(0,1)$ の座標を $s$ とする．両端近傍で零で $\int_0^1f(s)\,ds=1$ となる
滑らかな関数を取り，さらに $B\cap C$ の二つの collars 上で $f=0$ として，
$\zeta=f(s)\,ds$ を $B$ 側へ零延長する．
これは node 近傍で零の closed form で，$\int_\beta\zeta=1$ である．
従って $\lambda_c=\lambda+c\zeta$ は同じ $d\lambda$ と node 上の同じ primitive germ を持つが，
恒等写像の primitive 差は $c\ne0$ なら exact ではない．
この変更は fibration 自体と constant-zero Taylor 不変量を保ち，残余作用の原点だけを $c$ 変える．
非零の $c$ について，この closed 補正を base の外周 collar でも零にはできない．
そこには $H_1(N)$ の generator を持つ regular torus があり，前補題に反するからである．
\end{example}

\begin{proposition}[全飽和近傍の移送と focus 不変量]
\label{ke:focus-whole-neighborhood-transport}
前補題の $N$ 上で $\omega=d\lambda$ とし，
$\Phi:N\to N'$ が $\Phi^*\omega'=\omega$ を満たす微分同相であるとする．
$N'$ の葉の分割を $F'=F\circ\Phi^{-1}$ と定める．
全 compact fibre，その Lagrangian strata，primitive vanishing cycle と monodromy は
$\Phi$ によって移される．向き，Eliasson 座標，base cut と return cycle も一緒に運べば，
正則化作用の smooth germ $S$ が保存される．
さらに $\omega'=d\lambda'$ なら，
\[
 \Phi^*\lambda'-\lambda\text{ が exact}
 \quad\Longleftrightarrow\quad
 \int_{\Phi\beta}\lambda'=\int_\beta\lambda.
\]
この場合，全 closed cycles の作用積分と残余作用の定数項も保存される．
\end{proposition}
\begin{proof}
各 fibre は厳密に $\Phi(F^{-1}(c))$ なので，その位相型・連結性・strata が移る．
適合する Hamiltonians とその flows も共役されるため，return periods が一致する．
規格化を明示すると，normal の向きを固定し，周期 $1$ の moment を $\mu(o)=0$ と正規化する．
Eliasson 基底座標 $(c_1,c_2)$ を $c_2=\mu/(2\pi)$ と取り，
$c=c_1+\sqrt{-1}c_2$，$\alpha$ を normal cycle とする．
$c_1$ は一般には元の residual submersion そのものではない．
同じ slit disk と logarithm の枝，return cycle $\beta_c$ を選び，
$A(c)=\int_{\beta_c}\lambda$ と置く．V\~u Ng\d{o}c の規約では
\begin{align}
 dA&=\tau_1\,dc_1+\tau_2\,dc_2,\notag\\
 dS&=(\tau_1+\log|c|)\,dc_1+(\tau_2-\arg c)\,dc_2,\notag\\
 S(c)&=A(c)-A(0)+\Re(c\operatorname{Log}c-c),\qquad S(0)=0.
 \label{ke:eq:focus-regularized-action}
\end{align}
$dS$ は原点を越えて滑らかな閉形式となる\cite[Proposition 3.1, Definition 3.1]{VuNgoc}．
Return periods と正則化の選択が一致するので $S'=S$ である．
独立の向きを揃えた Eliasson 座標を使う際の flat な自由度については，
同文献の Lemma~4.1 と \S4 により Taylor series の一致を得る．

$\zeta=\Phi^*\lambda'-\lambda$ は closed なので，前補題が exactness の同値を与える．
一般には $A'(c)-A(c)=\int_\beta\zeta$ は定数であり，
$A'(0)-A(0)$ も同じ値なので $S$ では相殺される．
Exact の場合は任意の closed cycle 上の積分差が零となり，$A(0)$ も保存される．
これは\cite[\S7, ``Exact version'']{VuNgoc}の追加定数項である．
\end{proof}

Return cycle の変更 $\beta\mapsto\beta+k\alpha$ は $A$ を $2\pi k c_2$，
logarithm の枝の変更 $\operatorname{Log}c\mapsto\operatorname{Log}c+2\pi\sqrt{-1}\ell$ は
正則化項を $-2\pi\ell c_2$ だけ変える．両方を共通に運び，
$\partial_{c_2}S(0)\in[0,2\pi)$ の規格化を揃える．
作用を $J=A/(2\pi)$ で測る場合には，正則化も $S/(2\pi)$ とする．
上の命題は既に構成した葉の分割の push-forward を扱い，
独立に先指定した別の focus fibration との同値を与えるものではない．
実際，\cite[Theorem 2.1, \S6]{VuNgoc}により，規格化した Taylor series が
$0$ と $X^2$ の模型を実現できる．両者は同じ局所 focus 型と一重 monodromy を持つが，
その半大域不変量は異なる．この既知の実現原理の具体例と，
例\ref{ke:focus-local-primitive-shift}の primitive 原点だけの変更は別の現象である．

定理\ref{ke:projected-source-neck}への適用では，source を含む閉 residual leaf の
annular 近傍に submersion $a(p)$ を取り，$F=(\mu,a(p)-a(q))$ とする．
Source が相異なる regular leaves 上にあることから，他の source と saddle を避ける
小さい base disk と，その全逆像 $N$ を取れる．
Annulus の両境界と normal cuts の像を避けて base disk を縮めると，
任意の compact base subset の逆像は元の compact $K$ 内で閉じるので $F|_N$ は proper である．
共通 Moser domain はこの $N$ の全開近傍を含み，
補題\ref{ke:circle-exact-transport}は
$(\eta,\mathcal A\lambda_t)$ から $(\omega_t,\lambda_t)$ への
$\phi^*\lambda_t=\mathcal A\lambda_t+dQ$ を与えるので，
命題\ref{ke:focus-whole-neighborhood-transport}を適用できる．
別の canonical volume frame で primitive を作り直す場合は，
系\ref{ke:transition-residual-period}の残余 winding を別に確認する．
以上は各固定 $t$ と固定 stage の移送に関する保存則であり，
異なる $t$ や stage の focus 不変量が一定であるとは主張しない．

\subsection{残る円の平均と同一 normal moment}

今度は $\rho_{1,s}$ を neck の領域 $U$ 全体に与えた normal 円作用とする．
End の開領域 $E\subset U$ は $\rho_1$ に飽和しており，
$E$ 上には $\rho_1$ と可換な自由な正則円作用 $\rho_2$ があり，
両者で自由な $T^2$ 作用を生成すると仮定する．
各平均と homotopy を $\mathcal A_j,K_j$ と書き，
\[
 \eta=\mathcal A_1\omega,\qquad
 \kappa=\mathcal A_2\eta=\mathcal A_1\mathcal A_2\omega,\qquad
 \lambda_1=\mathcal A_1\lambda,\qquad
 \lambda_{12}=\mathcal A_2\lambda_1
\]
と置く．$\rho_2$ は $E$ のみに必要である．

第一の Moser flow $\phi_s$ は補題\ref{ke:circle-exact-transport}を
normal action に適用して $U$ 内に構成する．
第二の flow $\psi_s$ は $E$ 内で
\begin{equation}\label{ke:eq:residual-moser-path}
 \eta_s=(1-s)\kappa+s\eta,\qquad
 \beta_2=K_2\eta=\mathcal A_1K_2\omega,\qquad
 \iota_{Z_s}\eta_s=\beta_2,\quad \psi_0=\mathrm{id}
\end{equation}
によって構成する．以下では $\psi_s$ の定義域を $T^2$ に飽和した内側領域に取り，
$\psi_1$ の像が $\phi_s$ の共通定義域に含まれる buffer を課す．
$\phi=\phi_1,\ \psi=\psi_1,\ \Theta=\phi\circ\psi$ と書く．

\begin{theorem}[共通 normal circle と全作用周期を保つ二段階平均]
\label{ke:staged-neck-averaging}
上記の仮定の下で，$\psi_s$ は $\rho_1$ と可換であり，
\[
 \Theta^*\omega=\kappa,\qquad
 \Theta^*\lambda=\lambda_{12}+dq
\]
となる一価関数 $q$ がある．元の各 $T^2$ 軌道 $\mathcal O$ に対して
$\Theta(\mathcal O)$ は actual $\omega$-Lagrangian torus であり，
neck 全体の同じ actual 円作用
$\widehat\rho_1=\phi\rho_1\phi^{-1}$ の軌道で飽和している．

Normal moment を
$\mu=\lambda_1(V_1)$，
$\widehat\mu=\mu\circ\phi^{-1}$ と規格化すると，
\begin{equation}\label{ke:eq:staged-moment-identity}
 \mu\circ\psi=\mathcal A_2\mu,\qquad
 \widehat\mu\circ\Theta=\mathcal A_2\mu.
\end{equation}
より一般に，整数 cycle $b=(b_1,b_2)\in\ZZ^2$，
$V_b=b_1V_1+b_2V_2$ に対し，
元の軌道上のその一周を $\gamma_{b,x}$ とすると
\begin{equation}\label{ke:eq:staged-period-identity}
 \int_{\Theta(\gamma_{b,x})}\lambda
 =\int_{\gamma_{b,x}}\lambda_{12}
 =\lambda_{12}(V_b)(x).
\end{equation}
従って各 end で action の加法定数を独立に取り直す必要はない．
\end{theorem}
\begin{proof}
$\eta,\kappa,\beta_2$ は $\rho_1$ に不変である．
非退化形式に関する vector field の一意性から $Z_s$ は $\rho_1$-invariant であり，
flow の一意性により $\psi_s$ は $\rho_1$ と可換である．
補題\ref{ke:circle-exact-transport}を各段に用いると
\[
 \phi^*\omega=\eta,\quad \psi^*\eta=\kappa,\qquad
 \phi^*\lambda=\lambda_1+dq_1,\quad
 \psi^*\lambda_1=\lambda_{12}+dq_2.
\]
したがって $q=q_2+q_1\circ\psi$ とすれば所要の exact identity となる．
$\lambda_{12}$ は全 $T^2$ に不変なので，可換な生成場に対して
\[
 \kappa(V_1,V_2)
 =V_1\lambda_{12}(V_2)-V_2\lambda_{12}(V_1)
     -\lambda_{12}([V_1,V_2])=0.
\]
実次元二の自由な軌道は $\kappa$-Lagrangian であり，
$\Theta^*\omega=\kappa$ によりその像も Lagrangian である．
さらに
\[
 \widehat\rho_{1,s}(\Theta x)
 =\phi\rho_{1,s}\psi x
 =\phi\psi\rho_{1,s}x
 =\Theta\rho_{1,s}x
\]
なので，同じ actual normal circle が各 torus に含まれる．

Moment の等式は加法定数まで直接計算する．
$\iota_{V_1}\eta=-d\mu$ と可換性から
\[
 \beta_2(V_1)=K_2(d\mu)=\mathcal A_2\mu-\mu.
\]
$\mu_s=(1-s)\mathcal A_2\mu+s\mu$ と置けば
$\iota_{V_1}\eta_s=-d\mu_s$ である．従って
\[
 \frac d{ds}(\mu_s\circ\psi_s)
 =\bigl(\mu-\mathcal A_2\mu+d\mu_s(Z_s)\bigr)\circ\psi_s
 =\bigl(\mu-\mathcal A_2\mu+\beta_2(V_1)\bigr)\circ\psi_s=0.
\]
初期値は $\mathcal A_2\mu$ なので
\eqref{ke:eq:staged-moment-identity}を得る．
最後に exact differential の閉 cycle 上の積分は零であり，
$\lambda_{12}(V_b)$ は軌道上で一定なので
\eqref{ke:eq:staged-period-identity}が従う．
\end{proof}

\begin{corollary}[同じ normal action を持つ開 overlap 上の一致]
\label{ke:staged-normal-overlap}
Neck と triple chart に同じ $\omega$ と同じ normal 円作用 $\rho_1$ を与え，
同じ homotopy kernel を用いるとする．
内側 overlap からの normal Moser trajectories と，
residual stage の像からの trajectories が
$\rho_1$ に飽和した共通の開 overlap に留まると仮定する．
このとき両側の normal flow $\phi_s$ と actual normal action は一致する．
Triple 側に定理\ref{ke:staged-neck-averaging}の tori を用いれば，
それらは neck 側の同じ actual normal circle を含む．
また primitive $\lambda$ と cycle markings も同じものを使えば，
moment と作用周期は\eqref{ke:eq:staged-moment-identity}，
\eqref{ke:eq:staged-period-identity}によって共通に規格化される．
独立に固定した neck foliation の第二の周期や fiber 自体の一致には，
さらに reduced leaves を同じものに選ぶ必要がある．
\end{corollary}
\begin{proof}
共通 overlap では $\mathcal A_1\omega$ と $K_1\omega$ が一致する．
従って normal Moser vector field も一致し，
trajectory がその overlap に留まるという仮定の下で ODE の一意性を使える．
残りは定理\ref{ke:staged-neck-averaging}による．
\end{proof}

\begin{lemma}[短い作用に対する二段階構成の変位]
\label{ke:staged-displacement}
Normal action に不変な Hermitian metric $h$ を $U$ に取り，
$E$ 上では全 $T^2$ に不変とする．
$ch\le g_\omega\le Ch$，
$|V_1|_h\le e_1$ on $U$，
$|V_2|_h\le e_2$ on $E$ と仮定する．
上の二段階 flow を収める内在的 $h$-buffer があるなら，
\begin{equation}\label{ke:eq:staged-displacement}
 d_h(x,\Theta x)\le\frac{C}{2c}(e_1+e_2).
\end{equation}
各元の torus が内側領域に含まれる場合，その像の ambient diameter は
\begin{equation}\label{ke:eq:staged-ambient-diameter}
 \operatorname{diam}_{g_\omega}\Theta(\mathcal O)
 \le\sqrt C(e_1+e_2)\left(\frac Cc+\frac12\right).
\end{equation}
Centered homotopies を用いれば各段の変位定数 $C/(2c)$ は $C/(4c)$ に改善できる．
\end{lemma}
\begin{proof}
平均と凸結合は同じ二側比較を満たす．
$|K_1\omega|_h\le Ce_1/2$，
$|K_2\eta|_h\le Ce_2/2$ なので，各段の vector field の長さは
$Ce_j/(2c)$ 以下である．二つの trajectory を順に結べば
\eqref{ke:eq:staged-displacement}となる．
元の同じ torus の二点は，各 circle を高々半周ずつ進むことで
$h$-長さ $(e_1+e_2)/2$ 以下の経路で結べる．
その経路の両端に Moser trajectories を付け，
$g_\omega\le Ch$ を用いると\eqref{ke:eq:staged-ambient-diameter}を得る．
Centered kernel の絶対値の積分は $1/4$ なので最後の主張も従う．
\end{proof}

\subsection{一段階平均との Hamiltonian 比較}

二段階の endpoint map は，一段階の Moser map と一般に同一ではない．
ただし exact primitive を共通に固定した場合には，次の比較ができる．

\begin{proposition}[平均の順序と整数基底に関する Hamiltonian 比較]
\label{ke:staged-hamiltonian-comparison}
$E$ 上の同じ $T^2$ 作用と $\omega=d\lambda$ を用いる．
一段階の homotopy $H=K_1+\mathcal A_1K_2$ による map を $\Phi$，
上記の二段階 map を $\Theta$ とする．
以下の path homotopy の全 flow と inverse compositions を収める共通 buffer があるなら，
共通のより小さい開領域上で $\Phi^{-1}\Theta$ は
$\kappa$ に関する Hamiltonian isotopy の endpoint である．
従って $\Theta\Phi^{-1}$ は $\omega$ に関する Hamiltonian map である．

同じ結論は，$T^2$ の任意の二つの固定した整数基底で作る一段階または
二段階の平均の間にも成立する．
Centered homotopies を用い，$\omega$ が全 $T^2$ に不変な飽和 collar がある場合，
比較 isotopy はその collar を点ごとに固定できる．
\end{proposition}
\begin{proof}
三角形 $0\le v\le u\le1$ 上で
\[
 \Omega_{u,v}=(1-u)\kappa+(u-v)\eta+v\omega,\qquad
 \Lambda_{u,v}=(1-u)\lambda_{12}+(u-v)\lambda_1+v\lambda
\]
を定める．係数が非負で和が $1$ なので $\Omega_{u,v}$ は K\"ahler であり，
$d\Lambda_{u,v}=\Omega_{u,v}$ である．
$(0,0)$ から $(1,1)$ への対角 path は一段階構成，
$(0,0)\to(1,0)\to(1,1)$ の二辺は二段階構成に対応する．
二辺は各区間の端で stationary になるよう滑らかに parameter 化する．
この path と対角 path の凸結合によって，両座標が単調な滑らかな path family
$(u_r(s),v_r(s))$ を得る．

$\beta_1=K_1\omega$，$\beta_2=K_2\eta$ とし，
各 path 上の Moser primitive を
\[
 B_{r,s}=\dot u_r(s)\beta_2+\dot v_r(s)\beta_1
\]
とする．Homotopy identity より
\begin{align*}
 B_{r,s}
 &=-\partial_s\Lambda_{u_r,v_r}
       -d\bigl(\dot u_r K_2\lambda_1+\dot v_rK_1\lambda\bigr),\\
 dB_{r,s}&=-\partial_s\Omega_{u_r,v_r}.
\end{align*}
補題\ref{ke:circle-exact-transport}と同じ微分計算により，
各 endpoint map $F_r$ は
\begin{equation}\label{ke:eq:staged-path-exactness}
 F_r^*\omega=\kappa,\qquad F_r^*\lambda=\lambda_{12}+dQ_r
\end{equation}
を満たし，$Q_r$ は $r$ に滑らかである．
両座標の総変化は各々 $1$ なので，補題\ref{ke:staged-displacement}の
総変位の上界を全 $r$ に共通に使える．
さらに $F_r(D_{\mathrm{small}})\subset F_0(D_{\mathrm{large}})$ を保証する
image buffer を取ることが，相対写像の定義に必要である．

$W_r=(\partial_rF_r)\circ F_r^{-1}$ と置く．
\eqref{ke:eq:staged-path-exactness}を $r$ で微分すると
\[
 F_r^*(\iota_{W_r}\omega)
 =d\bigl(\partial_rQ_r-F_r^*(\lambda(W_r))\bigr).
\]
従ってこの deformation は Hamiltonian であり，
$F_0^{-1}F_r$ は $\kappa$-Hamiltonian isotopy となる．
端点は $F_0=\Phi$，$F_1=\Theta$ である．

整数基底を変えても全 Haar average $\mathcal A=\mathcal A_1\mathcal A_2$ は同じである．
二つの homotopies $H,H'$ について
\[
 H\omega=\mathcal A\lambda-\lambda-d(H\lambda),\qquad
 H'\omega=\mathcal A\lambda-\lambda-d(H'\lambda)
\]
なので，その差は exact である．
同じ線形 path $(1-s)\kappa+s\omega$ に対して primitive を
$(1-r)H\omega+rH'\omega$ と補間すれば，
再び endpoint の exact identity と上の Hamiltonian 計算が成立する．
二段階の場合は既に証明した比較を前後に組み合わせればよい．

最後に，centered homotopy は invariant form を零に写す．
指定 collar 上では $\beta_1=\beta_2=H\omega=H'\omega=0$ なので，
すべての vector fields と endpoint maps はそこで零および恒等となる．
従って比較 isotopy も同 collar を点ごとに固定する．
\end{proof}

\subsection{指定した normal subgroup を持つ actual closed corner}

\begin{corollary}[二段階 closed-corner bundle と同じ GH 誤差]
\label{ke:staged-closed-corner-gh}
系\ref{ke:actual-closed-corner}の固定次数 $d\ge5$ の族と exact triple chart を用いる．
その標準 $T^2$ の任意の固定した primitive normal subgroup $\alpha$ を選び，
$\alpha,\beta$ が $\ZZ^2$ の整数基底になるよう $\beta$ を補う．
$L=-\log|t|^2$，$R=\log L$ とし，
第一段を $\alpha$，第二段を $\beta$ とする上記の二段階 map $\Theta_t$ を取る．
十分小さい $t\ne0$ についてこの map は既存の softmax domain
$\mathcal H_t=\mathcal U_t(H_{L,R})$ の開近傍に定義できる．
\[
 \mathcal W_{t,J}^{\mathrm{st}}=\Theta_t(\mathcal H_t),\qquad
 \pi_{t,J}^{\mathrm{st}}=P_t\circ\Theta_t^{-1}
\]
とすると，その全 fibers は actual Lagrangian tori であり，
同じ actual normal action
$\widehat\rho_\alpha=\phi_t^\alpha\rho_\alpha(\phi_t^\alpha)^{-1}$ を含む．

模型 $\gamma_t=g_L^{\mathrm{sf}}$ に対し
\begin{equation}\label{ke:eq:staged-actual-metric}
 (1-C/R)\gamma_t\le\Theta_t^*g_t\le(1+C/R)\gamma_t.
\end{equation}
$R^{-2}g_t$ による $\mathcal W_{t,J}^{\mathrm{st}}$ の内在距離，
および $\mathcal V_t=\{\min\ell_a>R/2\}$ 内の全経路を許す距離のどちらについても
\begin{equation}\label{ke:eq:staged-closed-distortion}
 \sup_{x,y\in\mathcal W_{t,J}^{\mathrm{st}}}
 \left|d_t(x,y)-d_J(\pi_{t,J}^{\mathrm{st}}x,\pi_{t,J}^{\mathrm{st}}y)\right|
 \le C\frac{\log R}{R}.
\end{equation}
Fiber 自身の非正規化 intrinsic diameter は $O(R^{-1})$，
正規化後は $O(R^{-2})$ である．
Raw map と retraction に関する
\eqref{ke:eq:actual-closed-corner-target}と同じ誤差評価も成立する．
定数には固定した整数基底への依存を許す．

この bundle は前節の bundle と同一の写像としては主張しない．
両者は三重点の共通の開近傍上で actual $\omega_t$-Hamiltonian map により比較でき，
その map は同じ softmax base label を持つ fibers を対応させる．
\end{corollary}
\begin{proof}
補題\ref{ke:negative-exact}により三重点の actual form は exact である．
$h=h_t$ を constrained cusp metric とする．
外側領域 $\min\ell_a>R/2$ での全深部比較と模型の一様幾何から
\[
 ch\le\gamma_t\le Ch,\qquad
 |\nabla^h\gamma_t|_h\le C,\qquad |g_t-\gamma_t|_h\le C/R,
\]
\[
 |\omega_t|_h+|\nabla^h\omega_t|_h\le C,\qquad
 |V_\alpha|_h+|\nabla^hV_\alpha|_h
 +|V_\beta|_h+|\nabla^hV_\beta|_h
 \le C/R.
\]
生成場の評価は\eqref{ke:eq:full-corner-generator-c1}の固定整数線形結合である．
$h$ は全 torus に不変なので，$\eta_t,\kappa_t$ と両 convex paths も
$h$ に関する同じ二側比較および一様 $C^1$ bounds を持つ．
Isometry の pullback は $\nabla^h$ と可換であり，
chain contraction の微分には
$|\nabla^hV|_h|\omega|_h+|V|_h|\nabla^h\omega|_h$ だけが現れる．
従って各段の primitive とその一階微分は $O(R^{-1})$ であり，
非退化な Moser 方程式とその共変微分から
\[
 |Y_{t,s}|_h+|\nabla^hY_{t,s}|_h
 \le C/R
\]
を得る．これは補題\ref{ke:open-cell-moser-c1}および
\eqref{ke:eq:full-corner-moser-c1}と同じ計算である．

Constrained cusp norm では各 $d\log\ell_a$ が一様有界である．
従って $\min\ell_a>R$ からの二つの trajectories は，
$R$ が大きければ $\min\ell_a>3R/4$ に留まり，
$\min\ell_a>R/2$ との間の固定幅の buffer により全時刻で存在する．
必要な inverse compositions と path homotopy にも，少し大きい内側領域を用いて
同じ buffer を与えられる．

各段で，$\nabla^h\gamma_t$ の上界と上記の $C^1$ 評価により
$|\mathcal L_Y\gamma_t|_{\gamma_t}\le C/R$ であるから，
Gronwall の評価と二つの flow の合成により
$\Theta_t^*\gamma_t=(1+O(R^{-1}))\gamma_t$ を得る．
Actual/model 比較を最後に引き戻せば
\eqref{ke:eq:staged-actual-metric}となる．
模型の全閉 corner の距離評価
\eqref{ke:eq:closed-corner-model-input}と一様な正規化 diameter を使えば，
内在距離について\eqref{ke:eq:staged-closed-distortion}が従う．

保護領域内の距離には，系\ref{ke:closed-corner-protected-gh}と同じ
clipped projection $\Psi$ を用いる．二つの trajectories の模型長さの和は $C/R$ 以下であり，
\eqref{ke:eq:protected-pruning-length}から
\[
 d_J(\Psi(\Theta_tz),\Psi(z))\le C/R^2
\]
となる．保護領域全体の actual/model 下比較と，
既に得た内在距離の上界を合わせれば同じ distortion bound となる．
Fiber intrinsic diameter は，模型の phase torus の直径 $C/R$ と
\eqref{ke:eq:staged-actual-metric}から従う．
Raw map の評価は補題\ref{ke:raw-values}と
補題\ref{ke:raw-lipschitz}を二つの短い trajectories に用いればよい．

最後に，一段階・二段階および整数基底の変更に対して
命題\ref{ke:staged-hamiltonian-comparison}を適用する．
共通のソース形式は全 Haar average $\mathcal A\omega_t$ であり，
同じ softmax torus の二つの像を対応させる actual Hamiltonian map を得る．
\end{proof}

\subsection{Reduced foliations による延長と actual 適用の条件}

\begin{corollary}[共通 moment strip 上の条件付き neck 延長]
\label{ke:staged-reduced-extension}
定理\ref{ke:staged-neck-averaging}の normal 平均後の
$(U,\eta,\rho_1,\mu)$ において，共通の compact moment interval $I$ を取り，
reduced surfaces が境界付き曲面 $P$ の滑らかな束として同定されているとする．
さらに $P$ に，境界 coorientations と内部の embedded marking を固定した
基準 marked Morse foliation $\mathcal F_0$ があり，
境界近傍では radial，内部では simple saddles のみを持ち，
一つの critical leaf に一つの saddle だけを含むと仮定する．
各境界成分の互いに素な開 collar に，二段階構成から得た reduced circle foliations があり，
基準の radial circles をそれらへ写す collar 同定を $I$ に滑らかに選べるとする．
この同定は基準の collar 標識と coorientations に適合し，
互いに素な少し広い collars に台を持つ $\Diff_0(P)$ の滑らかな族へ延長できると仮定する．
このとき，それらの開 collar germs と一致する marked simple-saddle foliation family を
$P$ 上へ延ばせる．その円持ち上げを $\phi$ で移すと，
actual $\omega$-Lagrangian neck fibration が得られ，end では上記の staged tori と一致する．
その regular fiber は $T^2$，simple-saddle fiber は $S^1\times\figeight$ である．
\end{corollary}
\begin{proof}
仮定した $\Diff_0(P)$ の族により $\mathcal F_0$ を移す．
得られる族は指定した開 end germs と一致し，
内部の marked decomposition と一 critical-leaf 一 saddle の性質を保つ．
補題\ref{neck:marked-interpolation}と同様に，parameter germs の追加の補間が必要なら
collar 標識を固定した isotopies を用いる．
各固定 $\mu$ の reduced surface では一次元葉は isotropic であり，
その normal circle lift は Lagrangian である
（補題\ref{neck:lift}の reduction の証明）．
最後に $\phi^*\omega=\eta$ を用いて actual form へ移す．
\end{proof}

\begin{remark}[Actual neck への適用条件]
\label{ke:staged-neck-scope}
正則 normal action と K\"ahler 性から，
moment は各 complexified normal orbit に沿って狭義単調になる．
しかし，それだけでは全 base points に共通な moment interval や，
退化 parameter に一様な flow buffer は得られない．
系\ref{ke:staged-reduced-extension}では，これらの geometric identifications を仮定している．
本研究の actual KE 族では，source を除いた固定 compact base 上の正則 normal model に対し，
単一の global canonical primitive と共通 moment interval は
関連 KE 主稿\cite[\texttt{pair:cor:global-primitive}，
\texttt{pair:cor:actual-common-moment}]{KELocalCurrent}で検証される．
前者は base の nowhere-vanishing holomorphic $1$-form $\tau$ に対する
$\Omega=d\log v\wedge\tau$ と
$\lambda_t=d^c\log(\omega_t^2/(\Omega\wedge\overline\Omega))$ を用い，
$(\mathcal A_1\lambda_t)(V_1)$ が canonical normal moment と一致する．
後者は固定した end radius の後で $t\to0$ とする評価である．
Actual 適用には，さらに neck 全体の normal Moser flow の生存と buffer，
triple の各 end における同一作用・reduced foliation の同時一致を用いる．
これらは補題\ref{ke:fixed-stage-source-neck}と
系\ref{ke:neck-corner-saturated-union}で検証する．
Critical actions の直線性を保つには，さらに節\ref{ng8:section}の
relative area と marking の条件が必要である．

系\ref{ke:staged-closed-corner-gh}は，三重点の標準 torus 内の任意の
指定 primitive subgroup に適合できることを示す．
別々に構成した neck action との開 collar 上の一致を，
一つの三重点の複数の end で同時に満たすことは，同系だけの結論ではない．
節\ref{ke:relative-corner-section}で平均順序を相対化し，
系\ref{ke:neck-corner-saturated-union}で全交わりを同じ葉に揃える．
その後の source・branch・core を含む全域接着は
定理\ref{ke:actual-global-fibration}による．
非定数 holomorphic unit による orbit partition の変更も，
単なる整数基底変更の Hamiltonian 比較では処理されない．
全曲面の距離と singular fibers の一様な ambient 径は，
定理\ref{ke:global-ambient-lower}，\ref{ke:global-ambient-upper}および
命題\ref{ke:collapsed-tree-full-fiber-collapse}で，別の解析入力を組み合わせて証明する．
\end{remark}
